% Section li: Li coefficients, moment pencils and positivity certificates
\providecommand{\liE}{\mathcal{E}}
\providecommand{\liD}{\mathcal{D}}
\providecommand{\liZ}{\mathcal{Z}}
\providecommand{\liT}{\mathbb{T}}
\providecommand{\lidist}[1]{\lVert #1\rVert_{\mathbb{T}}}
\providecommand{\liEx}{\mathbb{E}}
\providecommand{\liind}{\mathbf{1}}
\providecommand{\liHk}{\mathcal{H}}
\providecommand{\liKk}{\mathcal{K}}

\section{Li coefficients, moment pencils and positivity certificates}\label{sec:li}

Li's criterion expresses the Riemann hypothesis as the positivity of a sequence of sums over
zeros \cite{Li1997Positivity,BombieriLagariasComplements,Lagarias2007Li}. This section collects
what sparse or finite portions of such sequences, and of the closely related power sums of
transformed zeros, decide. \Cref{li:sec:sampling,li:sec:curv} determine the index sets on
which an off-line zero must produce negative Li coefficients, and the sign patterns that the
actual coefficients of $\zeta$ repeat. \Cref{li:sec:ritz,li:sec:match} recover individual
zeros from Li moments: a moment pencil converges to each zero at a sharp rate, finite
inequalities certify simplicity, and mixed moments of two functions certify the separation
of their ordinates. \Cref{li:sec:hermite,li:sec:cluster} turn near-positivity of power-sum
matrices into quantitative zero location and complete zero counts, and show what finite data
cannot decide. \Cref{li:sec:univ} shows that the Cayley angles of the zeros generate a L\'evy
process with a universal small-time law. The recovery of a Dirichlet character from filtered
Li coefficients is in \cref{sec:rec}.

Most results below are statements about an abstract spectrum and are unconditional. They
become statements about $\zeta$ and Dirichlet $L$-functions through \cref{src:thm:dgrh},
which places the zeros on the critical line; such statements are marked \stG.

\paragraph{Notation.} For a zero $\rho$ put $u_\rho=1-1/\rho$. Then
\begin{equation}\label{li:eq:cayley}
 \abs{u_\rho}^2=1+\frac{1-2\Re\rho}{\abs\rho^2},\qquad \abs{1-u_\rho}=\abs\rho^{-1},
\end{equation}
so the critical line is mapped onto the unit circle. Write
$\xi(s)=\frac12s(s-1)\pi^{-s/2}\Gamma(s/2)\zeta(s)$. The Li coefficients are defined by
$\log[\xi(1/(1-z))/\xi(1)]=\sum_{n\ge1}\lambda_nz^n/n$, equivalently
$\lambda_n=\frac1{(n-1)!}\bigl(\frac{\dd}{\dd s}\bigr)^n\bigl[s^{n-1}\log\xi(s)\bigr]_{s=1}$;
for a number field $K$ the coefficients $\lambda_n(K)$ are defined in the same way from the
completed Dedekind zeta function $\xi_K$. Euler's constant is $\gamma$, or $\gamma_E$ where
$\gamma$ denotes an ordinate, and $\gamma_1,\gamma_2$ are the Stieltjes constants in
$\zeta(1+x)=x^{-1}+\gamma-\gamma_1x+\frac12\gamma_2x^2+O(x^3)$.

\subsection{Sampling Li coefficients}\label{li:sec:sampling}
\sources{grh_method_transfers_research; li_sampling_and_curvature}

\begin{definition}[Symmetric strip spectra]\label{li:def:spectrum}
A \emph{symmetric strip spectrum} is a locally finite multiset $\liZ$ in the open strip
$0<\Re\rho<1$, with finite multiplicities, invariant under $\rho\mapsto\bar\rho$ and
$\rho\mapsto1-\rho$, such that
\begin{equation}\label{li:eq:spectrum}
 W=\sum_{\rho\in\liZ}\abs\rho^{-2}<\infty .
\end{equation}
Sums over $\liZ$ count multiplicity. Its Li coefficients are
\begin{equation}\label{li:eq:Ln}
 L_0=0,\qquad L_n=\sum_{\rho\in\liZ}\bigl(1-\Re u_\rho^n\bigr)\quad(n\ge1),\qquad L_{-n}=L_n .
\end{equation}
\end{definition}

For Riemann's $\xi$ these are the classical $\lambda_n$. For a number field $K$ the
nontrivial zeros of $\zeta_K$ form a symmetric strip spectrum, and $L_n$ is the Li
coefficient $\lambda_n(K)=\lim_{T\to\infty}\sum_{\abs{\Im\rho}\le T}(1-u_\rho^n)$, which is
real. For a nonreal primitive Dirichlet character the union of the nontrivial zero multisets
of $L(s,\chi)$ and $L(s,\bar\chi)$ has both symmetries. The usual $O(T\log(T+2))$ zero count
for a fixed completed zeta or Dirichlet $L$-function gives \eqref{li:eq:spectrum} by partial
summation \cite{Lagarias2007Li}.

\begin{lemma}[Dominant phases; \stU]\label{li:lem:dom}
If every $\rho\in\liZ$ lies on the critical line, then $L_n\ge0$ for all $n$. Otherwise there
is a largest radius $R>1$ of the points $u_\rho$. List the arguments of the points on the
circle of radius $R$, with multiplicity, as $2\pi\alpha_1,\dots,2\pi\alpha_m$. There is
$1<r<R$ with
\[
 R^{-n}L_n=-\sum_{j=1}^m\cos(2\pi n\alpha_j)+O_{\liZ}\bigl(n^2(r/R)^n\bigr).
\]
\end{lemma}

\begin{proof}
On the line every summand of \eqref{li:eq:Ln} is nonnegative. Otherwise
\eqref{li:eq:cayley} and the reflection $\rho\mapsto1-\rho$ give a point with $\abs{u_\rho}>1$.
Since $\abs{1-u_\rho}=\abs\rho^{-1}$, only finitely many $u_\rho$ lie outside any fixed
neighbourhood of $1$, by \eqref{li:eq:spectrum}. This gives a finite dominant circle and a
strict gap to the remaining radii. For the remaining points the segment from $1$ to $u_\rho$
lies in $\abs z\le r$, and Taylor's integral remainder gives
\[
 \bigl|u_\rho^n-1+n/\rho\bigr|\le\tfrac12n(n-1)r^{n-2}\abs\rho^{-2}\qquad(n\ge2).
\]
Also $0<\Re(1/\rho)<\abs\rho^{-2}$. Thus \eqref{li:eq:Ln} converges absolutely, and summation
over the complement of the dominant circle gives the stated error.
\end{proof}

\begin{definition}[Schedules and Bohr recurrence]\label{li:def:bohr}
A \emph{schedule} is a family of positive indices $n_1(k),\dots,n_s(k)$, $s\ge1$, with
$\min_in_i(k)\to\infty$; dropping finitely many parameters $k$ gives the same eventual test.
Write $\liT=\R/\Z$ and $\lidist x=\min_{j\in\Z}\abs{x-j}$. The schedule has
\emph{simultaneous Bohr recurrence} if for every $d\ge1$, $\alpha\in\liT^d$ and $\delta>0$
there are arbitrarily large $k$ with
\begin{equation}\label{li:eq:return}
 \max_{i\le s,\,j\le d}\lidist{n_i(k)\alpha_j}<\delta .
\end{equation}
\end{definition}

\begin{theorem}[Universal Li sampling; \stU]\label{li:thm:sampling}
For every schedule the following are equivalent.
\begin{enumerate}[label=\textup{(\roman*)}]
\item The schedule has simultaneous Bohr recurrence.
\item For every symmetric strip spectrum with an off-line zero, $\max_iL_{n_i(k)}<0$ for
 arbitrarily large $k$.
\item For every such spectrum, with $R,m$ as in \cref{li:lem:dom},
 \[
  \liminf_{k\to\infty}\ \max_{1\le i\le s}R^{-n_i(k)}L_{n_i(k)}=-m .
 \]
\end{enumerate}
When \textup{(i)} fails there is a single finite off-line spectrum for which at least one
sampled coefficient is strictly positive at every parameter $k$. Its zeros may be taken
arbitrarily close to the critical line.
\end{theorem}

\begin{proof}
Apply \eqref{li:eq:return} to the dominant phases of \cref{li:lem:dom} with tolerances tending
to zero. All the cosines then approach one at every coordinate of the schedule. Their sum is
at most $m$, which gives (iii). Since $m\ge1$, (iii) implies (ii).

Suppose (i) fails. Choose phases $\alpha_1,\dots,\alpha_d$ and $0<\delta<1/2$ for which
\eqref{li:eq:return} eventually fails. If necessary append one irrational phase and shrink
$\delta$ to exclude the finitely many exceptional parameters: all their indices are positive,
so none of their multiples of that phase is zero. Then for every $k$ some $n=n_i(k)$ has
$n\alpha$ outside the open $\delta$-cube about zero. Consider
\[
 K_b(x)=\prod_{j=1}^d\Bigl(\frac{1+\cos(2\pi x_j)}2\Bigr)^b=\sum_ha_he^{2\pi ih\cdot x}\qquad(b\in\N).
\]
All $a_h$ are nonnegative dyadic rationals and $a_{-h}=a_h$. In one coordinate the constant
coefficient $4^{-b}\binom{2b}b$ is the largest, by the binomial theorem, and there are $2b+1$
coefficients of total mass one; hence $a_0\ge(2b+1)^{-d}$. Outside the cube
$K_b(x)\le\cos^{2b}(\pi\delta)$. Choose $b$ so that this is strictly below $a_0$. Put
$c=\sum_{h:\,h\cdot\alpha\in\Z}a_h\ge a_0$. Clearing a common denominator $D_0$ and collecting
equal phases gives
\[
 T(n)=D_0\bigl(K_b(n\alpha)-c\bigr)=\sum_{v\in V}b_vv^n,\qquad\abs v=1,\quad v\ne1,\quad
 b_v=b_{\bar v}\in\N .
\]
The finite set $V$ is nonempty, since $T(n)<0$ at every missed index just described.

Take the Cayley points $Rv$ and $R^{-1}v$, each with multiplicity $b_v$, where $R>1$ is
sufficiently close to one. At $R=1$ each inverse image $1/(1-v)$ has real part $1/2$.
Continuity and finiteness put all the perturbed inverse images inside the strip, as close to
the line as desired. Conjugation and inversion preserve the Cayley multiset, so both zero
symmetries hold. With $B=\sum_vb_v$, direct substitution gives
\begin{equation}\label{li:eq:counter}
 L_n=2B-(R^n+R^{-n})T(n).
\end{equation}
At each $k$ at least one sampled value exceeds $2B$. This contradicts (ii).
\end{proof}

\begin{corollary}[Finite perturbations of Riemann's $\xi$; \stG]\label{li:cor:xi}
If a schedule fails \textup{(i)} of \cref{li:thm:sampling}, a counterexample can be taken to be
an entire function of order one, real on the real axis and invariant under $z\mapsto1-z$,
whose zeros agree with those of $\xi$, with multiplicities, above some finite height, and whose
number of zeros in $\abs{\Im z}\le T$ equals that of $\xi$ for all sufficiently large $T$.
\end{corollary}

\begin{proof}
Let $A$ be the monic polynomial with the $2B$ zeros of \eqref{li:eq:counter}. Choose $B$
conjugate pairs of nonreal zeros of $\xi$, respecting their multiplicities, and let $H$ be
their monic polynomial. Both polynomials are real, monic of degree $2B$ and even about $1/2$.
The function $F=\xi A/H$ is entire with the stated symmetry and zero counts. At infinity
$A/H=1+O(\abs z^{-2})$, since the leading terms and the sums of the roots agree, so the order is
preserved. The retained zeros of $\xi$ lie on the line by \cref{src:thm:dgrh} and contribute
nonnegatively at every index, so the strictly positive witness survives.
\end{proof}

The construction uses the functional equation and the zeros only; an Euler product would be
an additional arithmetic constraint (\cref{li:q:euler}). \Cref{li:thm:surgery} gives a
related obstruction for the multiplicity of zeros.

\paragraph{A prescribed sparse index set.} Put
\begin{equation}\label{li:eq:E}
 M_j=2^{2^{j^2}},\qquad\liE=\bigcup_{j\ge1}\{M_j,2M_j,\dots,jM_j\}.
\end{equation}
The set $\liE$ contains $O(1+\log\log N)$ indices up to $N$: $M_j\le N$ implies
$j^2\le\log_2\log_2N$, and the first $j$ blocks contain $j(j+1)/2$ indices. It is the index set
of \cref{ana:thm:blocks}, whose extremal values of the prime remainder detect an off-line zero
only when the largest real part of a zero is attained. Li coefficients on the same set detect
every failure, and do so on a positive proportion of each late block.

\begin{theorem}[Sparse Li criterion; \stU]\label{li:thm:sparse}
For every number field $K$ the following are equivalent.
\begin{enumerate}[label=\textup{(\roman*)}]
\item All nontrivial zeros of $\zeta_K$ lie on the critical line.
\item $\lambda_n(K)\ge0$ for all sufficiently large $n\in\liE$.
\item The proportion of indices $k\in\{1,\dots,j\}$ with $\lambda_{kM_j}(K)<0$ tends to zero as
 $j\to\infty$.
\end{enumerate}
If \textup{(i)} fails, let $\varrho>1$ be the largest value of $\abs{u_\rho}$ and $m$ the number
of zeros attaining it, with multiplicity. Then
\[
 \lim_{j\to\infty}\min_{1\le k\le j}\varrho^{-kM_j}\lambda_{kM_j}(K)=-m .
\]
\end{theorem}

\begin{proof}
By \cref{li:lem:dom}, (i) implies (ii) and (iii). Suppose (i) fails. By \cref{li:lem:dom},
$\varrho^{-n}\lambda_n(K)=-\sum_{i=1}^m\cos(2\pi n\alpha_i)+o(1)$, which is at least $-m-o(1)$.
Let $\theta_1,\dots,\theta_t$ be the distinct values among the $\alpha_i$ and take an integer
$q\ge5$. The fractional parts of $0,M_j\theta,2M_j\theta,\dots,jM_j\theta$ are $N=j+1$ points in
$q^t$ half-open boxes of side $1/q$ in $[0,1)^t$. By the Cauchy--Schwarz inequality the number
of pairs in a common box is at least $(N^2/q^t-N)/2$. Each positive difference of the
multipliers occurs in at most $N$ pairs, so at least $(N/q^t-1)/2$ values $1\le k\le j$ satisfy
$\lidist{kM_j\theta_i}\le1/q$ for all $i$. For these $k$ the normalized coefficient is at most
$-m\cos(2\pi/q)+o(1)<0$. Letting $q\to\infty$ gives the limit $-m$, and the positive proportion
$(2q^t)^{-1}$ of negative coefficients in every late block contradicts (ii) and (iii).
\end{proof}

The proof allows any sequence of positive integer dilations $M_j\to\infty$ with block lengths
tending to infinity; the same pigeonhole argument shows that $\liE$, enumerated increasingly,
has simultaneous Bohr recurrence. For a field such as $\Q(\sqrt[3]2)$ the sampled
inequalities in (ii) are open.

\paragraph{Polynomial schedules.} For $E\subset\N$ the lower Banach density is
\[
 \underline d_{\mathrm B}(E)=\liminf_{H\to\infty}\inf_{a\in\Z_{\ge0}}\frac{\#(E\cap(a,a+H])}H .
\]
For sets of primes use the lower relative density
$\liminf_{X\to\infty}\#\{p\in E:p\le X\}/\pi(X)$.
These measure the abundance of good parameters; the sampled indices themselves, such as the
squares, may have density zero.

\begin{theorem}[The common-root criterion; \stU]\label{li:thm:poly}
Let $P_1,\dots,P_s\in\Z[x]$ be nonconstant with positive leading coefficients, and ignore
parameters at which some index is nonpositive.
\begin{enumerate}[label=\textup{(\alph*)}]
\item Over integers $k\ge1$, the equivalent conditions of \cref{li:thm:sampling} hold for
 $n_i(k)=P_i(k)$ if and only if the $P_i$ have a common root modulo every $q\ge1$.
\item Over primes $p$, they hold if and only if the $P_i$ have a common root coprime to $q$
 modulo every $q\ge1$.
\end{enumerate}
Under the corresponding condition, for every off-line spectrum and $0<\eta<1$ the set where
$\max_iR^{-P_i(k)}L_{P_i(k)}<-(1-\eta)m$ has positive lower Banach density in \textup{(a)} and
positive lower relative density among the primes in \textup{(b)}. One set of good parameters
works for any finite family of off-line spectra, each with its own $R,m$.
\end{theorem}

\begin{proof}
For necessity take the phase $1/q$ in \eqref{li:eq:return} with tolerance $1/(2q)$; a prime
larger than $q$ is a unit modulo $q$. \Cref{li:thm:sampling} supplies a strip counterexample if
the test fails.

For sufficiency fix $\alpha\in\liT^d$. Choose real representatives and a rational basis
$1,\beta_1,\dots,\beta_t$ of their rational span together with $1$. Write
$D\alpha_j=b_j+\sum_\ell a_{j\ell}\beta_\ell$ with integral coefficients and $D\ge1$. Choose a
common root $u\bmod D$, a unit in the prime case. In that class consider
$X(k)=(P_i(k)\alpha_j)_{ij}$. For a character $h=(h_{ij})\in\Z^{sd}$ put
$C_\ell(x)=\sum_{i,j}h_{ij}a_{j\ell}P_i(x)$. If $C_\ell$ is constant, evaluation at a common
root modulo each $q$ shows that the constant is zero. Set
$\Lambda=\{h:C_\ell\equiv0\text{ for all }\ell\}$ and
$H=\{x\in\liT^{sd}:h\cdot x=0\bmod1\text{ for all }h\in\Lambda\}$. On $k\equiv u\bmod D$ the
rational term $D^{-1}\sum h_{ij}b_jP_i(k)$ is integral. If $h\notin\Lambda$, the polynomial
$D^{-1}\sum_\ell\beta_\ell C_\ell(x)$ has an irrational leading coefficient in positive degree,
by rational independence. Along $k=u+Dt'$ the corresponding character has mean zero by Weyl's
theorem \cite{Weyl1916Gleichverteilung}, uniformly over the starting point of a long interval in
$t'$: in degree one the geometric sum is bounded uniformly, and in higher degree induction by the
van der Corput difference inequality applies, because for each fixed nonzero step the difference
has an irrational leading coefficient independent of the starting point. Taking the interval
length to infinity first, then the number of differences, gives uniform convergence.

For primes in the same reduced class, expand the indicator of $p\equiv u\bmod D$ as a sum of
additive characters in $p$. Each added term $ap/D$ preserves the irrationality of a nonconstant
coefficient. Rhin's theorem \cite{Rhin1973Repartition} gives mean zero over primes for each
resulting exponential, and the prime number theorem in arithmetic progressions gives the class
relative mass $1/\varphi(D)$. Characters in $\Lambda$ are identically one on the class.

In both settings the limiting probability measure on the progression is Haar probability on
$H$: its Fourier transform is one on $\Lambda$ and zero off $\Lambda$, by the annihilator identity
for subgroups of $\Z^{sd}$. Approximation by trigonometric polynomials gives convergence for
continuous functions, uniformly over starting points in the integer case; this includes $t=0$,
when all phases are rational. Every relative neighbourhood of zero in $H$ has positive measure.
Testing against a continuous nonnegative function supported in the chosen cube and positive at
zero gives positive lower density of returns on the progression, hence in the full parameter
set. For the last assertion concatenate all dominant angles and apply \cref{li:lem:dom}.
\end{proof}

\begin{example}[Schedules that detect and schedules that fail]\label{li:ex:schedules}
Detection means simultaneously negative values for every off-line spectrum.
\begin{center}\small
\begin{tabular}{@{}lll@{}}
\toprule
Indices at parameter $x$ & Integer parameters & Prime parameters\\
\midrule
$x^2,\ x^3$ & detect & fail modulo $2$\\
$(x-1)^2,\ x^2-1$ & detect & detect\\
$x,\ x+1$ & fail modulo $2$ & fail modulo $2$\\
$J(x),\ (x+2)J(x)$ & detect & detect\\
\bottomrule
\end{tabular}
\end{center}
A prime can use the residue class $0\bmod2$ only once. Here
$J(x)=(x^2-13)(x^2-17)(x^2-221)$ has no integer root but a unit root modulo every $q$. At an odd
prime $\ell\nmid221$ one of $13,17,221$ is a nonzero square modulo $\ell$, since if the first two
are nonsquares their product is a square; the root of the relevant quadratic is simple and lifts
to every power of $\ell$. At $13$ use $17\equiv4$, at $17$ use $13\equiv8^2$, and at $2$ use
$17\equiv1\bmod8$. The Chinese remainder theorem gives a unit root modulo every $q$.

For a failed modulus $q\ge2$ take Cayley points of both radii $R,R^{-1}$ above all nonidentity
$q$th roots of unity, with $R>1$ rational and close to one. The phase sum is $q-1$ at multiples of
$q$ and $-1$ elsewhere, so
\begin{equation}\label{li:eq:cyclo}
 L_n=\begin{cases}(q-1)(2-R^n-R^{-n})<0,&q\mid n,\\ 2(q-1)+R^n+R^{-n}>0,&q\nmid n.\end{cases}
\end{equation}
The zeros are the roots of a rational polynomial proportional to
$\prod_{t\in\{R,R^{-1}\}}\frac{(z-1)^q-t^qz^q}{(z-1)-tz}$. For $q=2$, $R=2$ this is $9z^2-9z+2$,
with zeros $1/3,2/3$ and Cayley points $-2,-1/2$; its Li coefficients are negative exactly at the
positive even indices, and no two consecutive ones are both negative.

Divisibility does not control arbitrary index sets. Put $v=(3+4i)/5=e^{2\pi i\theta}$, of
infinite order since $v+\bar v=6/5$ is not an algebraic integer. The set
$S=\{n:\Re(v^n)\le-1/2\}$ meets every $q\N$ infinitely often, yet the Cayley points
$R^{\pm1}v,R^{\pm1}\bar v$ with $R=6/5$ give $L_n=4-2(R^n+R^{-n})\Re(v^n)>0$ on $S$. The four
zeros form the rectangle $\{7/25,18/25\}+\{24i/25,-24i/25\}$.
\end{example}

\paragraph{Arithmetic equivalence from sparse windows.} For two number fields $K$ and $L$ put
$D_n=\lambda_n(K)-\lambda_n(L)$. The next result needs no hypothesis on the location of zeros.

\begin{theorem}[Arithmetic equivalence from sparse Li coefficients; \stU]\label{li:thm:rigidity}
If $\zeta_K\ne\zeta_L$, then
\[
 \lim_{H\to\infty}\ \inf_{A\ge1}\ \frac1H\sum_{n=A}^{A+H-1}\abs{D_n}=+\infty\qquad(A,H\in\N).
\]
Thus a bounded mean absolute discrepancy on a single sequence of windows of increasing length
forces $\zeta_K=\zeta_L$, and then $D_n=0$ for every $n$. The windows may be arbitrarily far
apart: those starting at $2^{2^{j^2}}$ and of length $j$ use only $O(1+\log\log N)$ indices up
to $N$. The same holds with $\abs{D_n}^2$ in place of $\abs{D_n}$.
\end{theorem}

\begin{proof}
The statement for squares follows from the one for absolute values by the Cauchy--Schwarz
inequality. Write $m_u$ for the difference of the multiplicities, for $K$ and for $L$, of the
transformed zero $u=1-1/\rho$. Second differences of the two symmetric sums give
\[
 a_n:=-(D_{n+2}-2D_{n+1}+D_n)=\sum_uc_uu^n,\qquad c_u=m_u(1-u)^2,\qquad\sum_u\abs{c_u}<\infty,
\]
because $\abs{1-u}=\abs\rho^{-1}$. All $\abs u$ are bounded. Suppose that there are windows
$[A,A+H-1]$ with $H\to\infty$ on which the mean of $\abs{D_n}$ is at most $C$. Then
$\abs{D_n}\le CH$ on each window.

\emph{Points off the unit circle.} Suppose some $u$ with $m_u\ne0$ has modulus above $1$, and
let $\varrho$ be the largest such modulus. It is attained at finitely many distinct points
$u_1,\dots,u_t$, and the other moduli with $m_u\ne0$ are at most some $r<\varrho$. Put
$z_i=u_i/\varrho$. For $N=A+H-t-2$ and $0\le h<t$,
\[
 \varrho^{-N-h}a_{N+h}=\sum_{i=1}^tz_i^h\,c_{u_i}z_i^N+O\bigl((r/\varrho)^{N+h}\bigr).
\]
The left sides tend to zero, since $N\ge H-t-1$ and $\abs{a_{N+h}}\le4CH$. The Vandermonde matrix
$(z_i^h)_{h,i}$ is invertible, so $c_{u_i}z_i^N\to0$, which contradicts $\abs{c_{u_i}}>0$. Hence
$m_u=0$ for $\abs u>1$. The reflection $\rho\mapsto1-\rho$ induces $u\mapsto1/u$ and preserves
multiplicity, so $m_u=0$ for $\abs u<1$ as well.

\emph{Points on the unit circle.} Let $\Gamma$ be the countable subgroup of the circle generated
by the remaining points $u$ with $m_u\ne0$, with the discrete topology. Its dual $G$ is compact
and metrizable. Write $g_n\in G$ for the character $u\mapsto u^n$. On each window take
nonnegative triangular weights $t_n$ of total mass one, obtained by normalizing
$\min(i+1,H-2-i)$, $0\le i\le H-3$, and zero at the last two positions. They satisfy
$\sup t_n=O(1/H)$ and $\sup\abs{t_{n+1}-t_n}=O(1/H^2)$. The measures $\sum_nt_nD_n\delta_{g_n}$
on $G$ have total variation $O(C)$, so a subsequence converges weak-star to a finite measure
$\nu$. For $z\in\Gamma$, Abel summation shows that the $t_n$-weighted average of $(u/z)^n$ tends
to zero for $u\ne z$, the total variation of the weights being $O(1/H)$. Since the $c_u$ are
absolutely summable, the weighted average of $a_nz^{-n}$ tends to $c_z$, with $c_z=0$ when
$m_z=0$. Shifting the indices in the second difference gives
\[
 \sum_nt_na_nz^{-n}=-(z-1)^2\sum_nt_nD_nz^{-n}+O\Bigl(H^{-2}\sum_n\abs{D_n}\Bigr),
\]
and the error is $O(C/H)$. Hence $\hat\nu(z)=-m_z$ for $z\ne1$. Subtracting $\hat\nu(1)$ times
the Haar probability makes the coefficient at $1$ vanish.

The resulting Fourier transform is bounded and integer-valued. Choose a polynomial that vanishes
at $0$ and takes the value $1$ at every nonzero integer of its finite range. Applying the
corresponding convolution polynomial to the measure gives a finite measure whose Fourier
transform is the indicator of $\{u:m_u\ne0\}$. By the idempotent theorem of Cohen, Helson and
Rudin \cite{GreenSanders2008Quantitative}, this set lies in the coset ring of $\Gamma$.

An infinite member of the coset ring of a subgroup of the circle is dense in the circle. To see
this, write a Boolean combination of finitely many cosets as a union of atoms. The positive
cosets of an atom intersect in a single coset of a subgroup $T$. Split it into finitely many
smaller cosets so as to settle all excluded cosets whose subgroups meet $T$ in finite index. Each
surviving piece is a coset of a subgroup $T_0$ with finitely many cosets of subgroups of infinite
index removed. If $T_0$ is infinite, it is dense in the circle. For a translation-invariant
finitely additive probability on the abelian group $T_0$, every coset of infinite index has
measure zero, by disjoint translates, and the trace on $T_0$ of every nonempty open arc has
positive measure, since finitely many translates cover $T_0$. So each surviving infinite piece is
dense. But the only possible accumulation point of $\{u:m_u\ne0\}$ is $1$, because
$u_\rho\to1$ as the height grows. Thus this set is finite.

\emph{Finitely many differing zeros.} The completed functions $\xi_K,\xi_L$ are entire of order
one, and Hadamard factorization now gives $\xi_K/\xi_L=e^{\alpha+\beta s}P(s)/Q(s)$ with
polynomials $P,Q$. Put $Z=\zeta_K/\zeta_L$. The gamma factors of a number field are products of
$\Gamma_\R(s)$ and $\Gamma_\R(s+1)$, and $\Gamma_\R(s+2)/\Gamma_\R(s)=s/(2\pi)$, while the
discriminant factor contributes a constant under $s\mapsto s+2$. Hence $Z(s+2)/Z(s)$ is a
rational function. On the positive real axis $\zeta_K(s)=1+O_K(2^{-s})$, and likewise for $L$, so
$Z(s+2)/Z(s)=1+O(2^{-s})$. A rational function with this decay towards $1$ is identically $1$.
For real $s>1$ this gives $Z(s)=Z(s+2m)\to1$, so $\zeta_K=\zeta_L$.
\end{proof}

The theorem identifies a zeta function up to equality, so it does not separate nonisomorphic
arithmetically equivalent fields. It concerns infinitely many samples and gives no effective
window length for a given pair of fields (\cref{li:q:rigidity}).

\subsection{Second differences of Li coefficients}\label{li:sec:curv}
\sources{li_sampling_and_curvature}

For a nonempty symmetric spectrum on the critical line write
\[
 d_n=L_{n+1}-2L_n+L_{n-1},\qquad C=d_0=2L_1,\qquad a_\rho=\abs{1-u_\rho}^2=\abs\rho^{-2}\qquad(n\in\Z).
\]
Centered second differences of \eqref{li:eq:Ln} give the absolutely and uniformly convergent
Fourier series
\begin{equation}\label{li:eq:fourier}
 d_n=\sum_\rho a_\rho u_\rho^n,\qquad C=\sum_\rho a_\rho ,
\end{equation}
which is real by conjugation. For $\xi$, $C=2\lambda_1=2+\gamma-\log(4\pi)$.

\begin{theorem}[Uniform recurrence of Li second differences; \stU]\label{li:thm:curv}
For every $h\in\Z$,
\begin{equation}\label{li:eq:curvcert}
 C-d_h\le\sup_{n\in\Z}\abs{d_{n+h}-d_n}\le\sqrt{2C(C-d_h)} .
\end{equation}
The upper bound is sharp on symmetric on-line spectra. For polynomial schedules satisfying
\cref{li:thm:poly}\textup{(a)} or \textup{(b)} and any $\epsilon>0$,
\[
 \max_i\sup_{n\in\Z}\abs{d_{n+P_i(k)}-d_n}<\epsilon
\]
on a set of positive lower Banach density, respectively of positive lower relative density
among the primes. The same parameters work for any prescribed finite family of such spectra. In
particular $\limsup_k\min_id_{P_i(k)}=C$.
\end{theorem}

\begin{proof}
Taking $n=0$ in the supremum gives the lower bound. For the upper bound apply the
Cauchy--Schwarz inequality to $\sum_\rho a_\rho\abs{u_\rho^h-1}$, using
$\sum_\rho a_\rho\abs{u_\rho^h-1}^2=2(C-d_h)$. A conjugate pair $u_\rho=e^{\pm2\pi i\theta}$ with
$\theta$ irrational has $d_n=C\cos(2\pi n\theta)$, and density of irrational rotations gives the
supremum $2C\abs{\sin(\pi h\theta)}$, which attains the bound. For recurrence truncate
\eqref{li:eq:fourier} so that the omitted mass is below $\epsilon/4$. Simultaneous returns make
$\sum_{\mathrm{retained}}a_\rho\abs{u_\rho^{P_i(k)}-1}<\epsilon/2$ for every $i$, and the tail costs
at most twice its mass, uniformly in $n$. Apply the density argument of the proof of
\cref{li:thm:poly} to the retained phases, and use their union for finitely many spectra.
\end{proof}

By \cref{src:thm:dgrh} this applies to $\xi$, to a real primitive Dirichlet character, and to the
paired zeros of a nonreal primitive character and its conjugate. For $\xi$ the ceiling $C$ is
approached arbitrarily closely, and $d_h<C$ for every $h\ne0$, since equality would require every
Cayley point to be an $h$th root of unity, whereas $\xi$ has infinitely many distinct zeros.
From now on in this subsection $d_n$ denotes the second differences of the Li coefficients of
$\xi$, and $\tau=C-d_1=4\lambda_1-\lambda_2$.

\begin{proposition}[Convexity windows from two Li coefficients; \stG]\label{li:prop:windows}
A set of integers $k$ of positive lower Banach density satisfies
\begin{align}
 d_{k^2+j},\ d_{k^3+j}&>0.00034&&(\abs j\le39),\label{li:eq:79window}\\
 d_{k^2+j},\ d_{k^3+j}&>0.00070&&(\abs j\le35).\label{li:eq:71window}
\end{align}
Every finite family of polynomials satisfying \cref{li:thm:poly}\textup{(a)} or \textup{(b)} may
replace $k^2,k^3$, for instance $(p-1)^2,p^2-1$ over primes.
\end{proposition}

\begin{proof}
Under RH the probability measure $\mu=C^{-1}\sum_\rho a_\rho\delta_{\Re u_\rho}$ on $[-1,1]$
satisfies $d_j/C=\int T_{\abs j}\,\dd\mu$, where $T_j(\cos\theta)=\cos(j\theta)$, and
$a:=\int x\,\dd\mu=d_1/C=1-\tau/C$. For $j\ge2$ put $b_j=\cos(\pi/j)$. The function equal to $-1$
on $[-1,b_j]$ and to $T_j$ on $[b_j,1]$ is a convex minorant of $T_j$: for $0<\theta<\pi/j$,
$T_j'(\cos\theta)=j\sin(j\theta)/\sin\theta=j\sum_{r=0}^{j-1}\cos((j-1-2r)\theta)$, each summand is
nonincreasing in $\theta$, so $T_j'$ is nondecreasing on $[b_j,1]$ and vanishes at $b_j$. Jensen's
inequality gives
\begin{equation}\label{li:eq:jensen}
 a\ge b_j\ \Longrightarrow\ d_j\ge CT_j(a).
\end{equation}
For $\abs u=1$, $\abs{1-u^j}\le\abs j\,\abs{1-u}$, and summation gives
$C-d_j=\frac12\sum_\rho a_\rho\abs{1-u_\rho^j}^2\le j^2\tau$.

Li's derivative formula gives $\lambda_1=L'(1)$ and $\lambda_2=2L'(1)+L''(1)$ for
$L=\log\xi$; expanding $(s-1)\zeta(s)$ at $s=1$ and using $\psi(1/2)=-\gamma-2\log2$ and
$\psi'(1/2)=\pi^2/2$ yields
\[
 \lambda_1=1+\frac{\gamma-\log(4\pi)}2,\qquad
 \lambda_2=2\lambda_1-1+\frac{\pi^2}8-\gamma^2-2\gamma_1 .
\]
Rational interval evaluation (Euler--Maclaurin at $N=256$ for $\gamma,\gamma_1$, Machin's
formula for $\pi$, the arctanh series for logarithms) certifies
\begin{gather*}
 0.0007430<C-35^2\tau<0.0007433,\qquad 0.0003500<CT_{39}(a)<0.0003504,\\
 -0.0015012<CT_{40}(a)<-0.0015008,
\end{gather*}
and $a>1-2/40^2\ge\cos(\pi/40)$, so $0\le j\arccos a<\pi$ for $0\le j\le40$. On $0\le j\le39$ the
least $T_j(a)$ occurs at $j=39$. This gives $d_j>0.00035$ for $\abs j\le39$ and $d_j>0.00074$ for
$\abs j\le35$; the value at $j=40$ shows where the two-moment bound stops certifying positivity.
\Cref{li:thm:curv} with uniform error below $10^{-5}$ transfers both windows to $k^2$ and $k^3$.
The certificate and its output are
\begin{center}
\artifact{li_sampling_and_curvature/constants_verify.py},\\
\artifact{li_sampling_and_curvature/constants_output.txt}.
\end{center}
\end{proof}

For $-1<a<1$ the bound \eqref{li:eq:jensen} is attained on $[b_j,1]$ by $\mu=\delta_a$, which is
realized by a single conjugate pair of critical-line zeros with Cayley phases $\pm\arccos a$.

\begin{proposition}[Negative second differences; \stG]\label{li:prop:negative}
One has $\liminf_{n\to\infty}d_n\le-\tau/\lambda_1<-0.00160638$. Moreover, for a set of $k$ of positive
lower Banach density, both windows of \cref{li:prop:windows} hold and
$d_{k^2+t},d_{k^3+t}<-0.00160$ for a fixed $t$.
\end{proposition}

\begin{proof}
Grouping equal phases in \eqref{li:eq:fourier} gives coefficients $c_u=m_u\abs{1-u}^2$, with $m_u$
the multiplicity. No phase is $1$, so the Ces\`aro mean of $d_n$ is zero, and the Ces\`aro mean
of $d_n^2$ is $V=\sum_uc_u^2\ge\sum_\rho a_\rho^2=2\tau$; finite geometric sums prove both mean
identities for partial sums, and absolute convergence passes them to the full series. Let
$B=-\liminf_nd_n$. For each $\varepsilon>0$ and all large $n$, $(C-d_n)(d_n+B+\varepsilon)\ge0$;
taking means and letting $\varepsilon\downarrow0$ gives $V\le CB$, hence $B\ge2\tau/C=\tau/\lambda_1$.
The numerical inequality is certified with the constants above. Fix $t$ with $d_t<-0.00160638$
and apply \cref{li:thm:curv} with uniform error below $10^{-6}$.
\end{proof}

\begin{theorem}[The first negative second difference; \stG]\label{li:thm:firstchange}
For the Riemann Li coefficients
\[
 d_j>0.00014\quad(0\le j\le117),\qquad d_{118}<-0.00026 .
\]
For a set of $k$ of positive lower Banach density, each of the two centers $k^2$ and $k^3$ lies
at the middle of a maximal run of exactly $235$ positive second differences, immediately preceded
and followed by negative values, with the same strict bounds. The same holds at the centers
$(p-1)^2$ and $(p-1)^3$ for a set of primes $p$ of positive lower relative density.
\end{theorem}

\begin{proof}
Let $t_1,\dots,t_{40}$ be the forty certified ordinates of \cref{li:cor:signs} below, and put
$w_j=2/(\frac14+t_j^2)$, $\theta_j=2\arctan(1/(2t_j))$. For a phase $\theta$ let $x=1-\cos\theta$.
Let $M_i$ be the total mass of $x^i$ in the measure of \eqref{li:eq:fourier}, and $r_i$ its mass
after deleting the forty pairs. The double-angle identity gives $M_0=C$, $M_1=C-d_1$,
$M_2=(3C-4d_1+d_2)/2$. With $L=\log\xi$ near $1$,
$L'''(1)=2-\frac74\zeta(3)+3\gamma_2+6\gamma\gamma_1+2\gamma^3$ and
$\lambda_3=3L'(1)+3L''(1)+\frac12L'''(1)$; this gives all three moments, in particular
$718389301443\cdot10^{-19}\le M_2\le718389301444\cdot10^{-19}$. For a pair, $x_j=w_j/4$, so
$r_i=M_i-\sum_{j\le40}w_j(w_j/4)^i$.

For $0\le x\le2$ and $n\ge0$ the Chebyshev polynomial satisfies
\begin{equation}\label{li:eq:chebtail}
 1-n^2x\le T_n(1-x)\le1-n^2x+\frac{n^2(n^2-1)}6x^2 .
\end{equation}
Indeed $\abs{T_n'}\le n^2$ on $[-1,1]$ follows from $\abs{\sin n\theta}\le n\abs{\sin\theta}$, and
$T_n'=n\sum_{j=0}^{n-1}T_{\abs{n-1-2j}}$ gives
$\max\abs{T_n''}\le n\sum_{j=0}^{n-1}(n-1-2j)^2=n^2(n^2-1)/3$; the mean value theorem and Taylor's
theorem at $1$ give the two bounds. Put $S_n=\sum_{j\le40}w_j\cos(n\theta_j)$. Summing
\eqref{li:eq:chebtail} against the omitted positive measure gives
\[
 S_n+r_0-n^2r_1\le d_n\le S_n+r_0-n^2r_1+\frac{n^2(n^2-1)}6r_2 .
\]
Rational interval evaluation at all $119$ indices gives the displayed inequalities, including
\[
 0.000147247112\le d_{117}\le0.000201258543,\qquad-0.000325756767\le d_{118}\le-0.000269874914 ;
\]
the program and the outward-rounded endpoints are
\begin{center}
\artifact{li_sampling_and_curvature/certify_window.py},\\
\artifact{li_sampling_and_curvature/window_bounds.csv}.
\end{center}

Since $d_{-j}=d_j$, the bilateral sequence has a run of $235$ positive terms on $[-117,117]$,
bracketed by two negative ones. Take a uniform recurrence tolerance below the least strict margin
in this pattern; \cref{li:thm:curv} with $k^2,k^3$ transfers every sign and bound to both
translated patterns. In the prime case the polynomials $(p-1)^2,(p-1)^3$ have the common unit
root $1$ modulo every integer.
\end{proof}

The finite sign pattern also follows, independently of the zero location, from the power series
of $\log(2\xi(s))$ at $s=-z/(1-z)$: its degree-$n$ coefficient is
$\sum_{j=1}^n\binom{n-1}{j-1}L^{(j)}(1)/j!=\lambda_n/n$. A $512$-bit ball-arithmetic evaluation
through degree $119$ gives, each within $5\cdot10^{-28}$,
\[
 d_{117}=0.000200146986527119643036209,\qquad d_{118}=-0.000271044467116152907821086,
\]
in agreement with the rigorous Keiper coefficients of Johansson \cite{Johansson2015Rigorous}. The
program is
\begin{center}
\artifact{li_sampling_and_curvature/verify_keiper.py}.
\end{center}

\begin{corollary}[Four sign patterns; \stG]\label{li:cor:signs}
For each $(a,b)\in\{-1,1\}^2$ the set $\{n\ge1:\ ad_n>11/1000,\ bd_{n^2}>11/1000\}$ has positive
lower Banach density. For the corresponding $h$ in the table below, the same inequalities hold
at $n=h+p-1$ for a set of primes $p$ of positive lower relative density, and one set of primes
serves all four rows.
\begin{center}\small
\begin{tabular}{rcc@{\qquad}rr}
\toprule
$h$ & $a$ & $b$ & lower bound for $ad_h$ & lower bound for $bd_{h^2}$\\
\midrule
891327 & $+$ & $+$ & 0.01161 & 0.01387\\
842728 & $+$ & $-$ & 0.01180 & 0.01349\\
602346 & $-$ & $+$ & 0.01165 & 0.01133\\
673696 & $-$ & $-$ & 0.01119 & 0.01148\\
\bottomrule
\end{tabular}
\end{center}
\end{corollary}

\begin{proof}
The polynomials $u$ and $2hu+u^2$ vanish at $u=0$, so \cref{li:thm:curv} gives a set of $u$ of
positive lower Banach density on which $d_{h+u}$ and $d_{h^2+2hu+u^2}$ are arbitrarily close to
$d_h$ and $d_{h^2}$; take the tolerance below the least margin of the row and put $n=h+u$. For
primes, $p-1$ and $(p-1)(2h+p-1)$ have the common unit root $1$ modulo every integer; taking the
family of $p-1$ and all four polynomials $(p-1)(2h+p-1)$ gives simultaneous returns for every row.

For the table, a pair of zeros with positive ordinate $t$ contributes $w(t)\cos(n\theta(t))$ to
$d_n$, with $w(t)=2/(\frac14+t^2)$ and $\theta(t)=2\arctan(1/(2t))$. For forty certified ordinates
$t_1,\dots,t_{40}$ put $W=C-\sum_jw(t_j)$; then for $e\in\{-1,1\}$
\begin{equation}\label{li:eq:signedtail}
 ed_n\ge\sum_{j=1}^{40}e\,w(t_j)\cos(n\theta(t_j))-W ,
\end{equation}
and $0.01022325<W<0.01022326$. The ordinates are certified without resolving all zeros below a
height. Forty centers $s_0=\frac12+it_0$ with $13<t_0<130$ receive disks of radius $10^{-18}$ with
disjoint ordinate intervals. At $s_0$ and $s_0+10^{-6}$, $\zeta$ is evaluated by Euler--Maclaurin
summation with cutoff $N=256$ and $20$ Bernoulli corrections, whose remainder is below $10^{-34}$,
in outward-rounded $192$-bit interval arithmetic. Every center has $\abs{\zeta(s_0)}<2\cdot10^{-25}$
and a difference quotient with a coordinate of modulus above $0.1$. First-order Euler summation
gives $\abs{\zeta(s)}<300$ on $\abs{s-s_0}\le1/4$, so by Cauchy's estimate the Taylor terms of
degree at least two are at most $5000\abs v^2$ for $\abs v\le10^{-6}$, and
$\abs{\zeta'(s_0)}>0.09$. On $\abs v=10^{-18}$ the linear term dominates, and Rouch\'e's theorem
gives exactly one zero in each disk. Inserting the ordinate intervals into
\eqref{li:eq:signedtail} gives the table
(\artifact{li_sampling_and_curvature/certify_zeta.py},
\artifact{li_sampling_and_curvature/certify_zeta_output.txt}).
\end{proof}

\subsection{Moment pencils and Ritz values}\label{li:sec:ritz}
\sources{li_ritz_zero_recovery; grh_new_results_2026-11-26; li_zero_separation}

Let $0<\gamma_1<\gamma_2<\cdots$ be the distinct positive ordinates of zeros of $\zeta$, and $m_i$
the multiplicity of $\frac12+i\gamma_i$. Under \cref{src:thm:dgrh} put
\begin{equation}\label{li:eq:mu}
 x_i=(\tfrac14+\gamma_i^2)^{-1},\qquad p_k=\sum_{i\ge1}m_ix_i^k\quad(k\ge1),\qquad
 \mu=\sum_{i\ge1}m_ix_i\,\delta_{x_i}.
\end{equation}
Thus $\int t^{k-1}\,\dd\mu=p_k$. The numbers $x_i$ are the denominators $\rho(1-\rho)$ of the
residues of the source transform at critical-line zeros. The Riemann--von Mangoldt count gives
$p_1<\infty$, and $\zeta$ has no real zero in $(0,1)$, as the alternating eta series shows.

\begin{proposition}[Moment coordinate; \stG]\label{li:prop:moments}
The function $E(z)=2\xi\bigl(\frac12+\sqrt{\frac14-z}\bigr)$ is entire and independent of the
choice of square root, and $E(z)=\prod_i(1-x_iz)^{m_i}$ with locally uniform convergence. Hence
\begin{equation}\label{li:eq:logmoments}
 p_k=-k[z^k]\log E(z)=\sum_{n=1}^k(-1)^{n+1}\binom{2k}{k-n}\lambda_n,\qquad
 p_1=\lambda_1=1+\frac\gamma2-\frac12\log(4\pi).
\end{equation}
\end{proposition}

\begin{proof}
The entire function $w\mapsto\xi(\frac12+w)$ is even, so its power series is a power series in
$w^2$, and substituting $w^2=\frac14-z$ shows that $E$ is entire. Under RH the paired Hadamard
product is $\xi(\frac12+w)=\xi(\frac12)\prod_{\gamma_n>0}(1+w^2/\gamma_n^2)$, with zeros listed with
multiplicity: pairing removes the exponential factors of genus one, evenness removes a linear
exponential prefactor, and $\sum\gamma_n^{-2}<\infty$; there is no central zero, since
$\zeta(1/2)\ne0$ (the alternating eta series at $1/2$ is positive while $1-\sqrt2<0$). Dividing by
the value at $w^2=\frac14$ and using $\xi(1)=\frac12$ gives
$2\xi(\frac12+w)=\prod_n(\gamma_n^2+w^2)/(\gamma_n^2+\frac14)$, which is the product for $E$.
Summability justifies the logarithmic series. The second formula in \eqref{li:eq:logmoments} is
Voros's inverse formula \cite[eq.~(11)]{Voros2006Sharpenings}: writing
$1-1/(\frac12+i\gamma_i)=e^{i\theta_i}$, one has $2-2\cos\theta_i=x_i$, the conjugate pair contributes
$2m_i(1-\cos n\theta_i)$ to $\lambda_n$, and the finite binomial expansion of $(2-2\cos\theta)^k$
gives the identity after summation; for fixed $n$ the pairing converges absolutely since
$1-\cos n\theta\le n^2(1-\cos\theta)$. Finally $p_1=-E'(0)/E(0)=\xi'(1)/\xi(1)=\lambda_1$, evaluated
with $\zeta(1+u)=u^{-1}+\gamma+O(u)$ and $\psi(1/2)=-\gamma-2\log2$.
\end{proof}

The alternating sum and monomial moment matrices both amplify roundoff; for computation the
Taylor coefficients of $\log E$ give the $p_k$ without subtracting rounded Li coefficients.

For $0\le a,b<N$ define the Hankel matrices $\liHk_N=(p_{a+b+1})$ and $\liKk_N=(p_{a+b+2})$, and
write $r_{1,N}>\cdots>r_{N,N}$ for the roots of $\det(\liKk_N-r\liHk_N)=0$. These are the
\emph{Ritz values} of multiplication by $t$ on polynomials of degree below $N$ in $L^2(\mu)$.

\begin{theorem}[Sharp error at every fixed rank; \stG]\label{li:thm:ritz}
For $N\ge j$ one has $0<r_{j,N}<x_j$, and $r_{j,N}$ increases to $x_j$. For every fixed $j$,
\begin{equation}\label{li:eq:ritzrate}
 \log(x_j-r_{j,N})=-4N\log N+4N\log\log N+O_j(N).
\end{equation}
The same expansion holds for $\log\bigl(\sqrt{1/r_{j,N}-1/4}-\gamma_j\bigr)$. Arbitrary finite
multiplicities are allowed.
\end{theorem}

Each pencil approximates several zeros simultaneously, and earlier zeros are not inputs. The
error law is for exact moments. The lower bound needs many distinct tail atoms at controlled
separation; counting zeros with multiplicity would not suffice.

\begin{lemma}\label{li:lem:packing}
For all large $T$ the interval $[T,2T]$ contains $\gg T\log T$ ordinates with consecutive selected
ordinates at least $1/\log(2T)$ apart.
\end{lemma}

\begin{proof}
Montgomery's Corollary~1, formula~(6) \cite{Montgomery1973Pair}, proved under RH, at height $U=2T$
and parameter $1/2$, bounds by $O(U\log U)$ the sum over ordered pairs $0<\gamma,\gamma'\le U$, with
multiplicity, of
\[
 \Bigl(\frac{\sin((\gamma-\gamma')\log U/4)}{(\gamma-\gamma')\log U/4}\Bigr)^2\frac4{4+(\gamma-\gamma')^2}.
\]
The summand is nonnegative and bounded below by a positive constant when
$\abs{\gamma-\gamma'}\le1/\log U$. Partition $[T,2T]$ into bins of that length, with a shorter final
bin. If the occupancies are $n_b$, then $\sum_bn_b^2\ll T\log T$, while $\sum_bn_b\asymp T\log T$ by
the Riemann--von Mangoldt formula; the Cauchy--Schwarz inequality gives $\gg T\log T$ occupied bins.
Keep the more numerous parity of the bin indices and one ordinate per retained bin.
\end{proof}

\begin{proof}[Proof of \cref{li:thm:ritz}]
\emph{Structure of the pencil.} On the space $V_N$ of polynomials of degree below $N$ with the
$L^2(\mu)$ inner product consider $R(P)=\int t\abs{P}^2\dd\mu/\int\abs P^2\dd\mu$. The pencil
represents the compression of multiplication by $t$ to $V_N$. Both Gram forms are positive definite,
since a nonzero polynomial cannot vanish on an infinite set. The full multiplication operator is
diagonal on the atoms with distinct eigenvalues $x_i$ and limit point zero, so min--max gives
$r_{j,N}\le x_j$ and $r_{j,N}\le r_{j,N+1}$. The monic orthogonal polynomial $Q_N$ of degree $N$
has the Ritz values as roots: $(t-r)P(t)$ for a Ritz vector $P$ is orthogonal to $V_N$, hence a
multiple of $Q_N$. This also gives one-dimensional eigenspaces and simple Ritz values, and the
three-term recurrence, with nonzero subdiagonal coefficients, gives strict interlacing for
consecutive $N$.

\emph{Upper estimate.} Write $w_i=m_ix_i$ and, for $i\le j$,
$P_{i,N}(t)=\prod_{k\le N,\,k\ne i}(t-x_k)/(x_i-x_k)$. For $0\le t\le x_{N+1}$ each numerator is
at most $x_k$, and for fixed $i$ each denominator is bounded below by a constant depending on $i$,
since $x_i$ is isolated. The zero count gives $x_k\ll(\log(k+1)/k)^2$ even for distinct ordinates
indexed without repetition. Hence, uniformly for $i\le j$ on this tail,
\begin{equation}\label{li:eq:ritzB}
 \abs{P_{i,N}(t)}\le B_{j,N},\qquad\log B_{j,N}\le-2N\log N+2N\log\log N+O_j(N),
\end{equation}
using $\sum_{k\le N}\log k=N\log N+O(N)$ and $\sum_{k\le N}\log\log(k+1)\le N\log\log(N+1)$. For
$P=\sum_{i\le j}a_iP_{i,N}$ the squared norm on the first $N$ atoms is $\sum_{i\le j}w_i\abs{a_i}^2$
and on the tail at most $p_1jB_{j,N}^2\sum_{i\le j}\abs{a_i}^2$. On the leading atoms where $P\ne0$
we have $t\ge x_j$. Hence $x_j-R(P)\le C_jB_{j,N}^2$ on this $j$-dimensional space, which gives the
upper bound for $x_j-r_{j,N}$ in \eqref{li:eq:ritzrate}.

\emph{Lower estimate.} Choose $T=DN/\log N$ with a large fixed $D$ and take $N$ ordinates from \cref{li:lem:packing}.
Their images $y_1,\dots,y_N$ under $t\mapsto(\frac14+t^2)^{-1}$, in monotone order, lie below
$x_{j+1}$ for large $N$ and satisfy, by the mean value theorem,
\begin{equation}\label{li:eq:ritzsep}
 y_i\gg(\log N/N)^2,\qquad\abs{y_i-y_k}\ge c\abs{i-k}\,d_N,\qquad d_N=\frac{(\log N)^2}{N^3}.
\end{equation}
For $P\in V_N$, interpolation at these nodes gives
\begin{equation}\label{li:eq:ritzinterp}
\begin{gathered}
 \abs{P(x_j)}\le\max_i\abs{P(y_i)}\sum_{i=1}^N\frac{(x_j/d_N)^{N-1}}{(i-1)!\,(N-i)!}
 =A_{j,N}\max_i\abs{P(y_i)},\\
 \log A_{j,N}\le2N\log N-2N\log\log N+O_j(N);
\end{gathered}
\end{equation}
the factorials contribute $2^{N-1}/(N-1)!$. Every $y_i$ carries mass at least $y_i$, so the squared
tail norm $S=\sum_{k>j}w_k\abs{P(x_k)}^2$ satisfies $S\ge(\min_iy_i)A_{j,N}^{-2}\abs{P(x_j)}^2$.
On the subspace of codimension $j-1$ where $P(x_1)=\cdots=P(x_{j-1})=0$,
\[
 x_j-R(P)\ge(x_j-x_{j+1})\frac S{w_j\abs{P(x_j)}^2+S},
\]
and the fraction is at least $c_N/(w_j+c_N)$ with $c_N=(\min_iy_i)A_{j,N}^{-2}$ (it is one if
$P(x_j)=0$). By the codimension form of min--max, $r_{j,N}$ is at most the maximal Rayleigh
quotient on this subspace, and \eqref{li:eq:ritzsep}, \eqref{li:eq:ritzinterp} give the lower bound
in \eqref{li:eq:ritzrate}. At finite $N$ that maximum is strictly below $x_j$, since $S>0$ for every
nonzero polynomial. Finally $x\mapsto\sqrt{1/x-1/4}$ has finite nonzero derivative at $x_j$.
\end{proof}

For fixed $N\ge j$ abbreviate $r_i=r_{i,N}$ and put
\begin{equation}\label{li:eq:septest}
 b_{j,N}=\min_{1\le k\le2N}\Bigl(\frac{p_k-\sum_{i<j}r_i^k}2\Bigr)^{1/k}.
\end{equation}
Form $S_j(t)=Q_N(t)/(t-r_j)$ and put $v_j=\int Q_N^2\,\dd\mu\big/\int S_j^2\,\dd\mu$. These are
exact functions of $p_1,\dots,p_{2N+1}$; for example
$\int Q_N^2\,\dd\mu=\det\liHk_{N+1}/\det\liHk_N$.

\begin{theorem}[Finite simplicity and enclosure test; \stG]\label{li:thm:ritzcert}
If $b_{j,N}<r_j$, then $m_1=\cdots=m_j=1$ and
\begin{equation}\label{li:eq:ritzcert}
 r_j\le x_j\le r_j+\frac{v_j}{r_j-b_{j,N}} .
\end{equation}
For fixed $j$ a passing test exists if and only if the first $j$ distinct ordinates are simple,
and then it passes for all sufficiently large $N$. The logarithm of the width in
\eqref{li:eq:ritzcert} has the same two leading terms and the same $O_j(N)$ precision as
\eqref{li:eq:ritzrate}.
\end{theorem}

\begin{proof}
Write $b=b_{j,N}$. In $p_k$ the first $j-1$ distinct sites contribute at least $\sum_{i<j}r_i^k$,
and the next two sites at least $2x_{j+1}^k$; thus $b\ge x_{j+1}$. If one of the top $j$
multiplicities exceeds one, then after one copy at each site $i<j$ is removed, at least two copies
at sites $\ge x_j$ remain, so $b\ge x_j\ge r_j$ and the test fails.

Under the probability measure $S_j^2\dd\mu/\int S_j^2\dd\mu$ the mean of $t$ is $r_j$, by
orthogonality of $Q_N$ to $S_j$, and the variance is $v_j$. When $b<r_j$ the support avoids
$(b,x_j)$, so $(t-b)(t-x_j)\ge0$ on it; taking expectations gives $v_j\ge(r_j-b)(x_j-r_j)$, which is
\eqref{li:eq:ritzcert}.

If the first $j$ sites are simple, $\sum_{i>j}m_i(x_i/x_j)^k\to0$ by dominated convergence. For a
fixed $k$ with this sum below one, the $k$th candidate in \eqref{li:eq:septest} tends to a value
strictly below $x_j$ as $N\to\infty$, so $r_j-b$ stays bounded away from zero. For the width, the
monic norm of $Q_N$ is minimal, and testing $\prod_{i\le N}(t-x_i)$ gives
$\int Q_N^2\dd\mu\le p_1\prod_{i\le N}x_i^2$. Also
$\int S_j^2\dd\mu\ge w_j\prod_{i\ne j,\,i\le N}\abs{x_j-r_i}^2\ge w_jc_j^{2(N-1)}$ eventually, since
the first $j-1$ Ritz values converge to points above $x_j$ and $r_i\le x_{j+1}$ for $i>j$. This
gives the upper rate; the lower rate follows from $x_j-r_j$ and \cref{li:thm:ritz}.
\end{proof}

The certificate uses that a multiplicity is at least one. An arbitrary positive measure can hide an
arbitrarily small extra atom above a proposed cutoff, so an ordinary moment support bound cannot
give the simplicity conclusion.

\begin{corollary}[Size of the Li moment determinant; \stG]\label{li:cor:hankel}
$\log\det\liHk_N=-2N^2\log N+2N^2\log\log N+O(N^2)$.
\end{corollary}

\begin{proof}
Let $h_n=\int Q_n^2\dd\mu$, with $Q_0=1$; then $\det\liHk_N=\prod_{n<N}h_n$. The trial polynomial
in the last proof bounds $h_n$ above. For the reverse bound use $n+1$ separated tail nodes from
\eqref{li:eq:ritzsep} with $N=n+1$. Comparing leading coefficients in the interpolation identity
for the monic $Q_n$ gives
$1=\sum_{i=1}^{n+1}Q_n(y_i)/\prod_{k\ne i}(y_i-y_k)$, so $\max_i\abs{Q_n(y_i)}\ge n!\,d_{n+1}^n/2^n$.
Each node has mass at least $\min_iy_i$. Hence $\log h_n=-4n\log n+4n\log\log n+O(n)$; sum over
$n<N$.
\end{proof}

\begin{corollary}[A local version for Dirichlet zeros; \stG]\label{li:cor:ritzdir}
Fix a primitive nonprincipal character $\chi$ of conductor $q$, a height $t_0\in\R$ and $a>1/2$.
Let $\gamma$ run over the ordinates of the completed function $\Lambda_\chi$, with multiplicity, and
assign to each distinct value $x=(a^2+(\gamma-t_0)^2)^{-1}$ the total multiplicity of its distance
$\abs{\gamma-t_0}$. The moments of this local distance spectrum are determined by Taylor data in
$\Re s>1$, and the $j$th Ritz value recovers the $j$th largest atom with
\[
 \lim_{N\to\infty}\frac{\log(x_j-r_{j,N})}{N\log N}=-4 .
\]
\Cref{li:thm:ritzcert} applies, with simplicity meaning total multiplicity one at a distance. All
assertions fix $q,t_0,a,j$ before the limit.
\end{corollary}

\begin{proof}
Put $G(w)=\Lambda_\chi(\frac12+it_0+w)\Lambda_{\bar\chi}(\frac12-it_0+w)$. The functional equation,
with the two root numbers multiplying to one, makes $G$ even. Then $E(z)=G(\sqrt{a^2-z})/G(a)$ is
entire of order at most $1/2$ with zeros $z=a^2+(\gamma-t_0)^2$; the two factors supply the two
symmetric $w$-roots of each zero, and a root at $w=0$ has twice the original multiplicity in $G$
and the original multiplicity in $E$. The genus-zero product and $E(0)=1$ give
$p_k=-k[z^k]\log E(z)$, computed from the completed functions at $\frac12+a\pm it_0$, where the
series converge absolutely. For the rate, the count of distances up to $T$ has main term
$(T/\pi)\log(qT/(2\pi e))$ with error $O_{q,t_0}(\log T)$, by sandwiching between centered counts
at $T\pm\abs{t_0}$ \cite{BennettMartinOBryantRechnitzer2021Counting}. This gives the upper bound
of \eqref{li:eq:ritzB}, and a distance in every interval of a fixed large length $L$ above a
threshold. Taking every other interval in a block of $2N$ such intervals starting at $N$ gives $N$
distinct distances of size $\asymp N$ separated by at least $L$, whose images have mass at least
$c/N^2$ and separation at least $c/N^3$. Interpolation gives $\log A_{j,N}\le2N\log N+O_j(N)$ and
the lower bound $-4N\log N-O_j(N)$.
\end{proof}

\begin{proposition}[A certificate from seventeen moments; \stG]\label{li:prop:ritz17}
The first four distinct positive ordinates of $\zeta$ are simple and satisfy
\begin{center}\small
\begin{tabular}{rrrr}
\toprule
$j$ & lower endpoint & upper endpoint & separating power $k$\\
\midrule
1 & 14.13472514173191 & 14.13472514173536 & 3\\
2 & 21.02201044341552 & 21.02204411662634 & 5\\
3 & 25.00089451502082 & 25.01282713231148 & 5\\
4 & 26.62873848826057 & 30.86048319178881 & 7\\
\bottomrule
\end{tabular}
\end{center}
\end{proposition}

\begin{proof}
The input is the first $17$ power sums of the scaled atoms $y=200x$. They come from the expansion
$\log(2\xi(1+u))=\sum_{k\ge1}(-1)^{k-1}p_k(u+u^2)^k/k$, with $Z(u)=u\zeta(1+u)$ expanded at $0$ by
Euler--Maclaurin with cutoff $n=256$ and last Bernoulli index $2p=192$. On $\abs u=1/2$ the
periodic Bernoulli remainder is at most $E=4(2p+1)!\,n/(6n)^{2p}$, because the normalized periodic
Bernoulli polynomial is bounded by $2\zeta(2p)/(2\pi)^{2p}<4/6^{2p}$, $\abs{(1+u)_{2p}}\le(2p+1)!$ and
the remaining integral is at most $n^{1-2p}$; the error of coefficient $k$ is at most $2^kE$.
Together with the expansions of $\log(1+u)$ and $\log\Gamma((1+u)/2)$, with integer zeta values
enclosed by rational Euler--Maclaurin bounds, all $17$ intervals for $200^kp_k$ have width below
$10^{-130}$.

For each $j$ a rational checker verifies three facts in the scaled variable. On a $j$-dimensional
rational trial space the form $\int(t-c_j)PQ\,\dd\mu$ is positive definite, by strict diagonal
dominance with positive diagonal; if at most $j-1$ atoms exceeded $c_j$, a nonzero combination
vanishing on them would make the form nonpositive, so the $j$th atom exceeds $c_j$. Next,
$p_k-\sum_{i<j}c_i^k\le2b^k$ for the listed $k$ and a rational $b<c_j$, which as in
\cref{li:thm:ritzcert} gives multiplicity one and an upper bound for the next atom. Finally the
variance about a rational trial center bounds the mean square deviation, and the product inequality
of \cref{li:thm:ritzcert} gives the upper bound. The data and checker are
\begin{center}
\artifact{li_ritz_zero_recovery/verification/certified_moments.json},\\
\artifact{li_ritz_zero_recovery/verification/witness.json},\\
\artifact{li_ritz_zero_recovery/verification/verify_zeros.py}.
\end{center}
\end{proof}

An uncertified floating-point run of the pencil with $N=8$ reproduces $\gamma_1$ to about $13$
digits from the upper side (\artifact{li_ritz_zero_recovery/verification/ritz_output.txt}).

Voros derived \eqref{li:eq:logmoments} and studied the associated secondary zeta functions
\cite{Voros2006Sharpenings,Voros2003Zeta}. Pavesi treats the same symmetric coordinate, Li
transforms and Hankel inertia \cite{Pavesi2026Klein}; Lasserre derives monotone support estimates
from generalized eigenvalues of moment matrices \cite{Lasserre2011Bounding}; Kawalec recovers
individual ordinates by limits of deflated power sums \cite{Kawalec2022Analytical}.

\subsection{Matched ordinates of two \texorpdfstring{$L$}{L}-functions}\label{li:sec:match}
\sources{li_zero_separation; li_ritz_zero_recovery}

Let $H$ be the entire completion of a primitive Dirichlet $L$-function, with the polar factors
removed in conductor one (so $H=\xi$ for $\zeta$). It has order one, its zeros are the nontrivial
zeros with multiplicity, and $H(s)=\epsilon\overline{H(1-\bar s)}$ with $\abs\epsilon=1$
\cite{Lagarias2007Li}. By \cref{src:thm:dgrh} each zero is $\rho=\frac12+i\gamma$ with $\gamma\in\R$.
For a center $c\in\R$ and a scale $a>0$ define
\[
 s_\pm=\tfrac12+ic\pm a,\quad d=2a,\quad S(u)=s_++\frac{du}{1-u},\quad
 z_\gamma=\frac{\rho-s_-}{\rho-s_+},\quad w_\gamma=\abs{1-z_\gamma}^2=\frac{d^2}{a^2+(\gamma-c)^2}.
\]
Then $\abs{z_\gamma}=1$, $0<w_\gamma\le4$ and $\sum_\gamma w_\gamma<\infty$; each occurrence of a zero
has its own label. Define the shifted Li coefficients
\[
 \ell_n=n[u^n]\log H(S(u))\ (n>0),\qquad\ell_0=0,\qquad\ell_{-n}=\overline{\ell_n},
\]
which do not depend on the branch of the logarithm. At $a=\frac12$, $c=0$ this is the usual Li
expansion; real shifts of the evaluation point go back to Sekatskii \cite{Sekatskii2014Generalized}.

\begin{lemma}[Circle moments of shifted Li coefficients; \stG]\label{li:lem:shifted}
Let $(Db)_n=2b_n-b_{n-1}-b_{n+1}$. For all $r\ge1$ and $n\in\Z$,
\begin{equation}\label{li:eq:shiftmoment}
 b_{r,n}:=-(D^r\ell)_n=\sum_\gamma w_\gamma^rz_\gamma^n .
\end{equation}
\end{lemma}

\begin{proof}
Write $G=H'/H$. The difference of genus-one logarithmic derivatives cancels the exponential
factor of the Hadamard product:
\[
 d\{G(S(u))-G(s_-)\}=d\sum_\rho\Bigl(\frac1{S(u)-\rho}-\frac1{s_--\rho}\Bigr)
 =\sum_\gamma\frac{w_\gamma}{1-z_\gamma u}\qquad(\abs u<1),
\]
since $S(u)-s_-=d/(1-u)$ and $w_\gamma=-d^2/[(\rho-s_-)(\rho-s_+)]$; both sums converge locally
uniformly. Differentiating the defining series gives $\sum_{n\ge1}\ell_nu^n=duG(S(u))/(1-u)^2$, so
for $n\ge1$ the coefficient of $u^n$ in $dG(S(u))$ is $\ell_{n+1}-2\ell_n+\ell_{n-1}$. At $u=0$ the
left side equals $2\Re\ell_1$, since $G(s_-)=-\overline{G(s_+)}$. Conjugation handles negative
indices, and $D$ multiplies the mode $z^n$ by $2-z-z^{-1}=\abs{1-z}^2$.
\end{proof}

For $r=1$, $a=\frac12$, $c=0$ and $H=\xi$, \eqref{li:eq:shiftmoment} is the Fourier series
\eqref{li:eq:fourier}. On the unit circle put $K_N(v)=\abs{N^{-1}\sum_{j=0}^{N-1}v^j}^2$ and
$K_N(v)^r=\sum_{\abs k\le r(N-1)}h_kv^k$; the rational coefficients $h_k$ are symmetric and
nonnegative, being the $r$-fold convolution of $(N-\abs k)/N^2$, and $K_N(1)=1$, $0\le K_N\le1$,
$K_N(v)\to0$ for $v\ne1$. For two primitive Dirichlet $L$-functions $A,B$ (possibly $\zeta$), with
centers $c_A,c_B$ and a common scale $a$, define the \emph{mixed test}
\begin{equation}\label{li:eq:F}
 F_{r,N}=\sum_{\abs k\le r(N-1)}h_kb^A_{r,k}\overline{b^B_{r,k}} .
\end{equation}
For the \emph{multiplicity test} take $A=B$, $c_A=c_B$, and subtract $b^A_{2r,0}$ from the right
side. Both use only indices $\abs n\le rN$.

\begin{lemma}[Pair expansion; \stG]\label{li:lem:pairs}
Write $z_i,v_j$ for the circle points of the zero labels of $A$ and $B$, $w_i,w'_j$ for their
weights, $p_{ij}=w_iw'_j$ and $q_{ij,N}=p_{ij}K_N(z_i\bar v_j)$. Then
$F_{r,N}=\sum_{i,j}q_{ij,N}^{\,r}\ge0$, the sum running over all ordered pairs in the mixed test
and over pairs of distinct labels in the multiplicity test.
\end{lemma}

\begin{proof}
Expand \eqref{li:eq:F} and apply \eqref{li:eq:shiftmoment}. All sums converge absolutely, since
$0\le q_{ij,N}\le p_{ij}\le16$ and $C:=\sum_{i,j}p_{ij}<\infty$. In the multiplicity test each
diagonal pair contributed $w_i^{2r}$; two distinct labels of the same zero survive.
\end{proof}

Put $x_i=\gamma_i-c_A$ and $y_j=\eta_j-c_B$. The map $x\mapsto(ix+a)/(ix-a)$ is injective, so
$z_i=v_j$ exactly when $x_i=y_j$, a \emph{match}. For $t>0$ only finitely many pairs have
$p_{ij}\ge t$. As $N$ grows, unequal pairs among them have $q_{ij,N}\to0$ and matched pairs keep
$q_{ij,N}=p_{ij}$; so for large $N$ the multiset of $q_{ij,N}\ge t$ is the multiset of matched
products $p_{ij}\ge t$.

\begin{theorem}[Completeness of the matched-pair test; \stG]\label{li:thm:match}
Let $T\ge0$ and $W=d^2/(a^2+T^2)$. In the mixed test there is no match with $\abs{x_i}\le T$ if and
only if $F_{r,N}<W^{2r}$ for some $r\ge1$, $N\ge2$. In the multiplicity test the same holds with
threshold $2W^{2r}$. On a clear window the strict inequality eventually holds along every joint
sequence $r,N\to\infty$.
\end{theorem}

\begin{proof}
A match in the window has $q_{ij,N}=w_i^2\ge W^2$, and a zero of multiplicity $m\ge2$ supplies
$m(m-1)\ge2$ ordered pairs; this gives necessity. Conversely let $R$ be the supremum of matched
products ($R=0$ if there is none); a positive supremum is attained. If the window is clear,
$R<W^2$, also when a match lies at the boundary. Choose $R<t<W^2$. Then $q_{ij,N}<t$ for all pairs
and all large $N$, and uniformly in these $N$,
$0\le F_{r,N}/W^{2r}\le(C/W^2)(t/W^2)^{r-1}\to0$.
\end{proof}

\begin{corollary}[Rate of the test; \stG]\label{li:cor:match}
Along every such joint sequence $F_{r,N}^{1/r}\to R$. If $R>0$ and $M$ ordered pairs realize the
largest matched product, there are $b<1$ and $N_0$ with $0\le F_{r,N}/R^r-M\le(C/R)b^{r-1}$ for
$N\ge N_0$. The least absolute centered ordinate of a match is $\sqrt{d^2/\sqrt R-a^2}$, read as
$+\infty$ when $R=0$.
\end{corollary}

\begin{proof}
For $R=0$ use $F_{r,N}\le Ct^{r-1}$ for any $t>0$. For $R>0$ remove the maximizing matched pairs;
the remaining matched supremum is strictly smaller, and the stabilization argument with threshold
$bR$ bounds the remainder. The formula inverts $R=w(x)^2$ with $w(x)=d^2/(a^2+x^2)$.
\end{proof}

Since $K_{mN}(v)=K_N(v)K_m(v^N)$, $F_{r,mN}\le F_{r,N}$, and for fixed $N$, $F_{r,N}^{1/r}$ is an
$\ell^r$ norm over pairs and decreases with $r$; so $N=2^j$, $r=j$ approach $R$ from above. No
uniform bound on the required $j$ is asserted.

\begin{theorem}[A gap bound from one strict test; \stG]\label{li:thm:gap}
Let $x_i,y_j\in[-T,T]$ be the centered ordinates of an admissible pair, and suppose a known
symmetry gives at least $k$ pair contributions equal to that of this pair in
\cref{li:lem:pairs} ($k=1$ always; $k=2$ in the multiplicity test, by reversal). Put $P_0=W^2$ and
$Q=(F_{r,N}/k)^{1/r}$. If $Q<P_0$, then
\begin{equation}\label{li:eq:gap}
 \abs{x_i-y_j}\ge d\sqrt{\frac{12}{N^2-1}\min\Bigl\{\frac{P_0-Q}{P_0^2},\frac{16-Q}{256}\Bigr\}} .
\end{equation}
\end{theorem}

\begin{proof}
Write $P=p_{ij}\in[P_0,16]$. For $v$ on the circle, $\abs{v^j-v^l}\le\abs{j-l}\,\abs{1-v}$, so
$1-K_N(v)=\frac1{2N^2}\sum_{j,l<N}\abs{v^j-v^l}^2\le\frac{N^2-1}{12}\abs{1-v}^2$. Subtracting the two
Cayley fractions gives $\abs{1-z_i\bar v_j}^2=\abs{z_i-v_j}^2=(P/d^2)\abs{x_i-y_j}^2$. By positivity
$Q\ge PK_N(z_i\bar v_j)$. Hence $\abs{x_i-y_j}^2\ge12d^2(P-Q)/[(N^2-1)P^2]$, and the derivative of
$(P-Q)/P^2$ has the sign of $2Q-P$, so its minimum on $[P_0,16]$ is at an endpoint.
\end{proof}

Let $\beta(s)=L(s,\chi_{-4})$, with completion
$H_\beta(s)=(4/\pi)^{(s+1)/2}\Gamma(\frac{s+1}2)4^{-s}\{\zeta(s,\frac14)-\zeta(s,\frac34)\}$. Both
$\xi$ and $H_\beta$ are real self-dual and centered at zero, and $\zeta(1/2)\ne0$, so simultaneous
negation gives $k=2$ in the mixed test for $\zeta,\beta$.

\begin{theorem}[Gaps between zeros of $\zeta$ and $\beta$; \stG]\label{li:thm:zetabeta}
If $\gamma$ and $\eta$ run over the ordinates of the nontrivial zeros of $\zeta$ and of $\beta$,
\[
 0.361<\min_{\abs\gamma,\abs\eta\le22}\abs{\gamma-\eta}<0.429 .
\]
Moreover all nontrivial zeros of $\zeta$ with $\abs\gamma\le30$ are simple, and any two of them with
distinct labels differ by more than $1.29$.
\end{theorem}

\begin{proof}
With $c_A=c_B=0$ and $k=2$, ball arithmetic evaluates the shifted Li coefficients by power-series
arithmetic, forms $b_{r,k}=-\sum_{j=-r}^r(-1)^j\binom{2r}{r+j}\ell_{k+j}$ and the integer kernel
polynomial, and certifies:
\begin{center}\small
\begin{tabular}{@{}lrrrrrrr@{}}
\toprule
functions & $T$ & $a$ & $N$ & $r$ & degree & $F/W^{2r}<$ & gap $>$\\
\midrule
$\zeta,\beta$ & 20 & 14.5 & 50 & 5 & 250 & 0.682658 & 0.477\\
$\zeta,\beta$ & 22 & 24 & 90 & 5 & 450 & 0.737047 & 0.361\\
$\zeta$, distinct labels & 30 & 24 & 30 & 3 & 90 & 1.151405 & 1.29\\
\bottomrule
\end{tabular}
\end{center}
The last column is \eqref{li:eq:gap}, rounded down. A multiple zero would give gap zero between
distinct labels in the last row. For the upper bound, the signs of $\xi$ at $\frac12+21.022i$ and
$\frac12+21.023i$ are $(-,+)$, and those of $H_\beta$ at $\frac12+21.450i$ and $\frac12+21.451i$ are
$(-,+)$; continuity gives zeros with $0<\eta-\gamma<0.429$. No ordinates enter the lower-bound
verifier. The program and its output are
\begin{center}
\artifact{li_zero_separation/verify_collisions.py},\\
\artifact{li_zero_separation/collisions-output.txt}.
\end{center}
\end{proof}

\begin{example}[A prescribed gap from three coefficients per center]\label{li:ex:shifted}
Take two copies of $\zeta$ with $c_A=21$, $c_B=14$, $a=T=0.15$, $N=3$, $r=k=1$. Taylor expansions
through degree three give $F/W^2<0.761429$, so every pair with $\abs{\gamma-21},\abs{\eta-14}\le0.15$
satisfies $\abs{(\gamma-\eta)-7}>0.082$. By \cref{li:prop:ritz17} the first two positive zeros of
$\zeta$ lie in these windows.
\end{example}

Brown used finite Li data for zero-free regions \cite{Brown2005Li}, and Ernvall-Hyt\"onen,
Od\v{z}ak, Smajlovi\'c and Su\v{s}i\'c studied Li coefficients of products of shifted functions
\cite{ErnvallHytonenOdzakSmajlovicSusic2015Modified}. A finite-height simplicity criterion for $\zeta$
by trace defects is in \cite{Alderliefste2026Finite}.

\subsection{Hermite power-sum matrices}\label{li:sec:hermite}
\sources{grh_zero_certificates_2026-11-26}

Let $Z=(z_1,\dots,z_n)$ be a multiset in $\C$, conjugation invariant in the sense that a point and
its conjugate have the same multiplicity. Put
\begin{equation}\label{li:eq:hermite}
\begin{gathered}
 s_k=\sum_{\ell=1}^nz_\ell^k,\quad s_0=n,\qquad H_n(Z)=(s_{i+j})_{0\le i,j<n},\\
 b(Z)=\max_\ell\abs{\Im z_\ell},\qquad\liD(Z)=\max\{0,-\lambda_{\min}(H_n(Z))\}.
\end{gathered}
\end{equation}
The $s_k$ are real, so $H_n(Z)$ is real symmetric. The defect $\liD$ is the least $e\ge0$ with
$H_n(Z)+eI$ positive semidefinite, and also the operator-norm distance to the positive
semidefinite cone.

\begin{theorem}[Two sharp powers; \stU]\label{li:thm:hermite}
Let $n\ge2$, let $Z$ be conjugation invariant with $\abs{z_\ell}\le1$ for all $\ell$, and put
$K_n=4\floor{n/2}$, $C_n=n^2(n-1)(7n-11)/6$. Then
\begin{equation}\label{li:eq:hmain}
 2\Bigl(\frac{b(Z)}{K_n}\Bigr)^{2n-2}\le\liD(Z)\le C_nb(Z)^2 ,
\end{equation}
and consequently
\begin{equation}\label{li:eq:hinverse}
 \sqrt{\liD(Z)/C_n}\le b(Z)\le\min\Bigl\{1,K_n\Bigl(\frac{\liD(Z)}2\Bigr)^{1/(2n-2)}\Bigr\}.
\end{equation}
Repeated roots are allowed and no separation is assumed. For each $n$ the exponent $1/(2n-2)$ on
the right of \eqref{li:eq:hinverse} and the exponent $2$ on the right of \eqref{li:eq:hmain}
cannot be increased. The constants are not asserted optimal.
\end{theorem}

In particular $\liD=0$ if and only if all roots are real; this is the classical Hermite--Sylvester
theorem, with Nathanson's interpolation proof \cite{Nathanson2021Hermite}. For $n=3$ the theorem
reads $b^4/128\le\liD\le30b^2$.

For a real polynomial $p(x)=\sum_{r<n}c_rx^r$,
\begin{equation}\label{li:eq:hquad}
 c^{\mathsf T}H_n(Z)c=\sum_{\ell=1}^np(z_\ell)^2 ,
\end{equation}
with ordinary squares: an interpolated value $i$ at a nonreal point contributes $-1$.

\begin{lemma}[Grouped divided differences]\label{li:lem:dd}
Assume $b=b(Z)>0$, let $d\le n$ be the number of distinct root locations, and list the distinct
nonnegative imaginary heights as $0=y_0<y_1<\cdots<y_k=b$, so $k\le\floor{d/2}$. Choose $j$ with
$\delta=y_j-y_{j-1}\ge b/k$, and put $L=\{z:\Im z\le-y_j\}$, $U=\{z:\Im z\ge y_j\}$,
$M=\{z:\abs{\Im z}\le y_{j-1}\}$. Prescribe the values $-i,0,i$ on $L,M,U$ and order the distinct
nodes with $L$ first, then $M$, then $U$. Every divided difference of positive order $r$ of these
data on a consecutive block satisfies $\abs{f[x_a,\dots,x_{a+r}]}\le2^{r-1}\delta^{-r}$, with no
separation of nodes inside a group.
\end{lemma}

\begin{proof}
The groups are exhaustive because no height lies in the selected gap; $L$ and $U$ are nonempty;
distinct groups are at distance at least $\delta$, and $L,U$ at distance at least $2\delta$. A
positive-order divided difference of constant data vanishes, so blocks inside a group contribute
nothing. In order one, neighbouring groups have value difference $1$ at distance at least $\delta$,
and $L,U$ have difference $2$ at distance at least $2\delta$. For a mixed block of order $r\ge2$ the
first and last nodes lie in different groups, at distance at least $\delta$, and the recursion with
the two bounds of order $r-1$ gives $2\cdot2^{r-2}\delta^{-(r-1)}/\delta$.
\end{proof}

\begin{proof}[Proof of \cref{li:thm:hermite}]
\emph{Lower bound.} If $b=0$, \eqref{li:eq:hquad} makes $H_n$ positive semidefinite. Otherwise
there is a unique polynomial $p$ of degree at most $d-1$ with the values of \cref{li:lem:dd}; the
data respect conjugation, so by uniqueness $p$ has real coefficients. Suppose $\abs{z_\ell}\le R$.
The coefficient norm of $\prod_{a<r}(x-x_a)$ is at most $(1+R)^r$, and Newton's expansion with
\cref{li:lem:dd} gives, for $T=2(1+R)/\delta>2$,
$\norm p_{\mathrm{coef},2}\le\norm p_{\mathrm{coef},1}\le1+\frac12\sum_{r=1}^{d-1}T^r\le T^{d-1}$.
By \eqref{li:eq:hquad} the padded coefficient vector has
$c^{\mathsf T}H_nc=-\#\{\ell:z_\ell\in L\cup U\}\le-2$, with multiplicity; interpolation used only
distinct locations. The Rayleigh quotient gives
\begin{equation}\label{li:eq:hradius}
 \liD(Z)\ge2\Bigl(\frac\delta{2(1+R)}\Bigr)^{2d-2}\ge2\Bigl(\frac b{2k(1+R)}\Bigr)^{2d-2},
\end{equation}
and replacing $k$ by $\floor{n/2}$, $d$ by $n$ and $R$ by $1$ gives the lower half.

\emph{Upper bound.} Put $a_\ell=\Re z_\ell$ and $v_r=\sum_\ell a_\ell^r$; the matrix $H_\R=(v_{i+j})$ is
positive semidefinite. For $r\ge2$ Taylor's formula on the vertical segment from $a_\ell$ to $z_\ell$,
inside the unit disk, gives $\abs{\Re(z_\ell^r)-a_\ell^r}\le\frac{r(r-1)}2\abs{\Im z_\ell}^2$, the linear
term being imaginary. As $s_r$ is real, $\abs{s_r-v_r}\le\frac{nr(r-1)}2b^2$, with equality of the
terms for $r\le1$. The maximal absolute row sum, attained on the last row, gives
$\norm{H_n-H_\R}_{\mathrm{op}}\le\frac{nb^2}2\sum_{j=0}^{n-1}(n-1+j)(n-2+j)=C_nb^2$, and positivity of
$H_\R$ gives $c^{\mathsf T}H_nc\ge-C_nb^2$ for unit $c$.
\end{proof}

\begin{proposition}[Sharpness of both exponents; \stU]\label{li:prop:hsharp}
The following examples show that both exponents of \cref{li:thm:hermite} are optimal.
\begin{enumerate}[label=\textup{(\roman*)}]
\item Fix $n\ge2$, distinct nonzero real numbers $a_1,\dots,a_{n-2}$, and a small $h>0$.
 Let $Z_t$ consist of $t$ times the roots $a_1,\dots,a_{n-2},ih,-ih$. Then $\liD(Z_t)\sim(-\sigma)t^{2n-2}$ as
 $t\downarrow0$, for a constant $\sigma<0$, while $b(Z_t)=ht$.
\item Adding $-\sigma t^{2n-2}$ to the single moment $s_{2n-2}$ of $Z_t$ makes the Hankel matrix
 positive semidefinite. So the exponent $1/(2n-2)$ is optimal even for the implication ``a measured
 Hankel matrix is positive semidefinite and every measured moment has error at most $e$''.
\item For $Z_b=\{ib,-ib,0,\dots,0\}$ with $n-2$ zeros and $0<b\le1$,
 $\liD(Z_b)=2b^2\sum_{r=0}^{\floor{n/2}-1}b^{4r}\sim2b^2$.
\end{enumerate}
\end{proposition}

\begin{proof}
(i) At $h=0$ the leading $(n-1)\times(n-1)$ block $A$ is positive definite: a nonzero real
polynomial of degree at most $n-2$ cannot vanish at the $n-1$ points $0,a_1,\dots,a_{n-2}$, which
carry positive weights. It stays positive definite for small $h$. Write
$H=\begin{pmatrix}A&v\\v^{\mathsf T}&c\end{pmatrix}$. The Vandermonde factorization shows $\det H<0$:
the conjugate pair has squared difference $-4h^2$, real nodes have positive squared differences,
and the factors linking a real node to the pair combine to $(a_j^2+h^2)^2>0$. So
$\sigma=c-v^{\mathsf T}A^{-1}v=\det H/\det A<0$; this includes $n=2$. Scaling gives
$H(t)=D_tHD_t$ with $D_t=\mathrm{diag}(1,t,\dots,t^{n-1})$, and with
$D_t'=\mathrm{diag}(1,\dots,t^{n-2})$ completion of the square gives
\[
 \binom xy^{\!\mathsf T}H(t)\binom xy=(D_t'x+t^{n-1}A^{-1}vy)^{\mathsf T}A(D_t'x+t^{n-1}A^{-1}vy)
 +\sigma t^{2n-2}y^2\ge\sigma t^{2n-2}(\norm x^2+y^2).
\]
The vector $y=1$, $x=-t^{n-1}(D_t')^{-1}A^{-1}v$ has $\norm x^2=O(t^2)$, so
$\sigma t^{2n-2}\le\lambda_{\min}(H(t))\le\sigma t^{2n-2}/(1+O(t^2))$. A bound $b\le C\liD^\alpha$ with
$\alpha>1/(2n-2)$ would give $t\ll t^{(2n-2)\alpha}$.

(ii) The change affects only the bottom-right entry and makes the Schur complement zero.

(iii) The odd power sums vanish, so $H_n$ splits into its even-index and odd-index blocks. With
$v=(1,-b^2,b^4,\dots)$ of the appropriate length, the even block is
$2vv^{\mathsf T}+(n-2)e_0e_0^{\mathsf T}$, positive semidefinite, and the odd block, of size
$\floor{n/2}$, is $-2b^2vv^{\mathsf T}$.
\end{proof}

\begin{example}[An exact cubic family]\label{li:ex:hcubic}
For the roots $\{t,it/2,-it/2\}$, $0<t\le1$,
\[
 H(t)=\begin{pmatrix}3&t&t^2/2\\t&t^2/2&t^3\\t^2/2&t^3&9t^4/8\end{pmatrix},
\]
the leading $2\times2$ block has determinant $t^2/2$ and the Schur complement is $-25t^4/8$. The
vector $u_t=(3t^2/2,-5t,1)^{\mathsf T}$ has $u_t^{\mathsf T}H(t)u_t=-25t^4/8$ and
$\norm{u_t}^2=1+25t^2+9t^4/4$, so
\[
 \frac{25t^4}{8(1+25t^2+9t^4/4)}\le\liD\le\frac{25t^4}8,\qquad b=t/2 .
\]
Replacing $9t^4/8$ by $17t^4/4$ gives a positive semidefinite matrix. The family $\{0,ib,-ib\}$
instead has $\liD=2b^2$. Exact checks are in
\begin{center}
\artifact{grh_zero_certificates_2026-11-26/data/exact_certificates.json}.
\end{center}
\end{example}

\begin{remark}[Fewer distinct locations and weights]\label{li:rem:distinct}
If at most $d\ge2$ distinct locations occur, \eqref{li:eq:hradius} gives
$\liD(Z)\ge2(b(Z)/(4\floor{d/2}))^{2d-2}$ even when $n>d$; repetition does not force the worst
exponent. For $d$ distinct conjugation-invariant nodes in the unit disk with positive
conjugation-invariant weights, minimum weight $w_*$ and total weight $W$, the $d\times d$ weighted
power-sum matrix satisfies $2w_*(b/(4\floor{d/2}))^{2d-2}\le\liD\le\frac{Wd(d-1)(7d-11)}6b^2$. For
roots in $\abs z\le R$ the lower constant becomes $2\floor{n/2}(1+R)$ and the upper constant
$\frac n2\max_{i<n}\sum_{j<n,\,i+j\ge2}(i+j)(i+j-1)R^{i+j-2}$. The defect depends on the monomial basis
and on the normalization of the disk.
\end{remark}

Suppose real approximate moments satisfy $\abs{s_k-\widehat s_k}\le e_k$ for $0\le k\le2n-2$, and put
$\widehat H=(\widehat s_{i+j})$ and
$E=\min\{\max_i\sum_je_{i+j},(\sum_{i,j}e_{i+j}^2)^{1/2}\}$, a bound for
$\norm{H-\widehat H}_{\mathrm{op}}$.

\begin{corollary}[Certificates from inaccurate moments; \stU]\label{li:cor:hnoise}
Under \cref{li:thm:hermite}, if a certified interval $[\ell,u]$ contains $\lambda_{\min}(\widehat H)$,
then with $D_-=\max\{0,-u-E\}$ and $D_+=\max\{0,-\ell+E\}$,
\begin{equation}\label{li:eq:hcertified}
 \sqrt{D_-/C_n}\le b(Z)\le\min\{1,K_n(D_+/2)^{1/(2n-2)}\}.
\end{equation}
A certified $\widehat H\succeq-\tau I$, $\tau\ge0$, gives
\begin{equation}\label{li:eq:hslack}
 b(Z)\le\min\{1,K_n((\tau+E)/2)^{1/(2n-2)}\},
\end{equation}
and $E=ne$ may be used if every $e_k\le e$. If $\widehat H\succeq\gamma I$ is certified with
$\gamma>E$, then all $n$ roots are real and distinct.
\end{corollary}

\begin{proof}
The quadratic forms of $H$ and $\widehat H$ differ by at most $E$ on unit vectors, so
$\lambda_{\min}(H)\in[\ell-E,u+E]$; apply \eqref{li:eq:hinverse}. If $\widehat H\succeq\gamma I$ with
$\gamma>E$, then $H\succ0$; the roots are real, and fewer than $n$ distinct locations would make $H$
singular by the Vandermonde factorization.
\end{proof}

For $n=3$, $\tau=0$ and uniform error $e$, \eqref{li:eq:hslack} gives $b\le4(3e/2)^{1/4}$, so
$e\le\eta^4/384$ certifies $b\le\eta$; by \cref{li:prop:hsharp}(ii) the fourth power cannot be
improved for this certificate. Symmetric Schur elimination decides positive semidefiniteness of a
rational matrix exactly.

\begin{theorem}[A local analytic certificate; \stU]\label{li:thm:contour}
Let $\Omega$ be a bounded conjugation-invariant domain in the unit disk with positively oriented
piecewise smooth boundary $\Gamma$. Let $F$ be holomorphic near $\overline\Omega$, with
$F(\bar z)=\overline{F(z)}$, nonzero on $\Gamma$, and with exactly $n\ge2$ zeros in $\Omega$ counted
with multiplicity. Then the moments
\[
 s_k=\frac1{2\pi i}\int_\Gamma z^k\frac{F'(z)}{F(z)}\,\dd z\qquad(0\le k\le2n-2)
\]
satisfy \cref{li:thm:hermite} and \cref{li:cor:hnoise}, with $b$ the largest absolute imaginary
part of a zero in $\Omega$.
\end{theorem}

\begin{proof}
The residues of $F'/F$ are the multiplicities, so the residue theorem gives \eqref{li:eq:hermite};
the zeros in $\Omega$ are conjugation invariant and lie in the unit disk.
\end{proof}

The zero count must be certified before choosing the matrix size; an error below $1/2$ in $s_0$
identifies it by rounding. For $n\le1$ conjugation symmetry already makes the zeros real. For
meromorphic $F$ poles would contribute negative weights. The contour representation and the
qualitative inertia test are due to Deng, Yang and L\"u \cite{DengYangLu2026Contour}.

\begin{corollary}[Horizontal distance of $\xi$ and Dirichlet zeros in a window; \stU]\label{li:cor:hstrip}
For a height $T\in\R$ and a scale $A>0$ let $F(z)=\Xi(T+Az)$ with $\Xi(w)=\xi(\frac12+iw)$, or
$F(z)=e^{-i\theta/2}\Lambda_\chi(\frac12+i(T+Az))$ for a primitive nonprincipal $\chi$ with root number
$e^{i\theta}$ and completion
$\Lambda_\chi(s)=(q/\pi)^{(s+\mathfrak a)/2}\Gamma((s+\mathfrak a)/2)L(s,\chi)$, $\chi(-1)=(-1)^{\mathfrak a}$.
Then $F$ is real entire. If its normalized contour moments over a domain as in
\cref{li:thm:contour} satisfy the certified slack condition \eqref{li:eq:hslack}, every enclosed zero
$\rho=\frac12+i(T+Az)$ satisfies
\[
 \abs{\Re\rho-\tfrac12}\le A\min\{1,K_n((\tau+E)/2)^{1/(2n-2)}\}.
\]
\end{corollary}

\begin{proof}
$\Xi$ is real entire by the functional equation and conjugation. For $\chi$,
$\overline{\Lambda_\chi(s)}=\Lambda_{\bar\chi}(\bar s)$ and $\Lambda_\chi(s)=e^{i\theta}\Lambda_{\bar\chi}(1-s)$
give $\overline{F(z)}=F(\bar z)$. A zero $z$ corresponds to $\rho$ with $\abs{\Re\rho-\frac12}=A\abs{\Im z}$.
\end{proof}

No hypothesis on zeros is used. By \cref{src:thm:dgrh} the left side is zero for these functions;
the finite statement applies to audited numerical data and to other real entire functions. The
phase cancels from $F'/F$. An imprimitive character should be reduced to its primitive completion,
since the extra Euler factors have their own zeros. The moments of high order are poorly conditioned
for large $n$, and smaller windows, each with certified boundary and count, give better bounds.

Certified Hermite matrices from approximate roots are due to Ayyildiz Akoglu and Szanto
\cite{AyyildizAkogluSzanto2023Certified}, and Hankel relaxations for zeros of entire functions under
a spacing condition to Bickel, Pascoe and Sargent \cite{BickelPascoeSargent2023Zero}.

\subsection{Complete cluster counts and what finite data cannot decide}\label{li:sec:cluster}
\sources{grh_new_results_2026-11-26}

Let $(x_n)$ be a finite or countable multiset of positive reals, with positive integer
multiplicities, such that
\begin{equation}\label{li:eq:summ}
 p_1=\sum_nx_n<\infty,\qquad p_j=\sum_nx_n^j\quad(j\ge1).
\end{equation}
Then $x_n\le p_1$ and only finitely many occurrences lie above any positive threshold. Choose
distinct centers $a_1,\dots,a_m>0$, radii $0<r_i<a_i$ with pairwise disjoint intervals
$I_i=(a_i-r_i,a_i+r_i)$, and an integer $K\ge1$. Put
\[
 P(x)=\prod_{i=1}^m(x-a_i),\quad D_i=\prod_{j\ne i}(a_i-a_j),\quad
 L_i(x)=\frac{\prod_{j\ne i}(x-a_j)}{D_i},\quad d_i=\frac K{a_i}+2\sum_{j\ne i}\frac1{a_i-a_j},
\]
\[
 Q_i(x)=\Bigl(\frac x{a_i}\Bigr)^K\bigl[1-d_i(x-a_i)\bigr]L_i(x)^2,\qquad
 R_K=\sum_nx_n^KP(x_n)^2\ge0,\qquad S_i=\sum_nQ_i(x_n).
\]
Then $Q_i(a_j)=\delta_{ij}$, $Q_i'(a_j)=0$ and $x^K\mid Q_i$, and $R_K$, $S_i$ use only the $2m+1$
moments $p_K,\dots,p_{K+2m}$. Put $b_i=a_i-r_i$, $g_{ij}=\abs{a_i-a_j}-r_i$,
$A_i=\prod_{j\ne i}g_{ij}^{-2}$, $T_i=K/b_i+2\sum_{j\ne i}g_{ij}^{-1}$,
$U_i=K/b_i^2+2\sum_{j\ne i}g_{ij}^{-2}$ and
\[
 C_i=\max\Bigl\{\frac{b_i^{-K}A_i}2(T_i^2+U_i),\ \frac{a_i^{-K}}{D_i^2}\bigl(r_i^{-2}+\abs{d_i}r_i^{-1}\bigr)\Bigr\}.
\]

\begin{theorem}[Residual certificate; \stU]\label{li:thm:residual}
Let $N_i$ be the number of occurrences in $I_i$. Then $\abs{N_i-S_i}\le C_iR_K$. For $\tau>0$ the
number $N_{\mathrm{miss}}$ of occurrences $x_n\ge\tau$ outside $\bigcup_iI_i$ satisfies
$N_{\mathrm{miss}}\le R_K/(\tau^K\prod_ir_i^2)$. Consequently, if $R_K<\tau^K\prod_ir_i^2$, the
intervals lie above $\tau$, and each interval $[S_i-C_iR_K,S_i+C_iR_K]$ contains exactly one
nonnegative integer $n_i$, then the spectrum at or above $\tau$ lies in the intervals with exact
counts $(n_1,\dots,n_m)$.
\end{theorem}

\begin{proof}
We show $\abs{\liind_{I_i}(x)-Q_i(x)}\le C_ix^KP(x)^2$ for $x>0$. Outside $I_i$, with $u=x-a_i$,
$\abs u\ge r_i$, and away from other centers,
$\abs{Q_i(x)}/(x^KP(x)^2)=a_i^{-K}D_i^{-2}\abs{1-d_iu}/u^2\le a_i^{-K}D_i^{-2}(r_i^{-2}+\abs{d_i}r_i^{-1})$;
at another center both sides vanish. Inside $I_i$ let $h_i(x)=x^{-K}L_i(x)^{-2}$, so
$h_i(a_i)=a_i^{-K}$ and $h_i'(a_i)=-d_ia_i^{-K}$; then $Q_i$ is $x^KL_i^2$ times the linear Taylor
polynomial of $h_i$ at $a_i$, while $1=x^KL_i^2h_i$. Since $P=(x-a_i)D_iL_i$, Taylor's theorem gives
$\abs{1-Q_i(x)}/(x^KP(x)^2)\le\frac1{2D_i^2}\sup_{I_i}\abs{h_i''}$, and logarithmic differentiation gives
\[
 \frac{h_i''(t)}{D_i^2}=t^{-K}\prod_{j\ne i}(t-a_j)^{-2}\Bigl\{\Bigl(\frac Kt+2\sum_{j\ne i}\frac1{t-a_j}\Bigr)^2
 +\frac K{t^2}+2\sum_{j\ne i}\frac1{(t-a_j)^2}\Bigr\}\le b_i^{-K}A_i(T_i^2+U_i).
\]
Summing the pointwise bound gives the count estimate. Each missed occurrence contributes at least
$\tau^K\prod_ir_i^2$ to $R_K$. A nonnegative integer below one is zero, and $N_i$ lies in its
certified interval.
\end{proof}

A count of one certifies one simple occurrence; a larger count may be a repeated value or several
nearby values. With $x^KP^2=\sum_ju_jx^j$ and $Q_i=\sum_jv_{ij}x^j$, and enclosures
$\abs{p_j-\widehat p_j}\le\epsilon_j$, put $\widehat R=\sum u_j\widehat p_j$,
$R^+=\widehat R+\sum\abs{u_j}\epsilon_j$, $\widehat S_i=\sum v_{ij}\widehat p_j$ and
$E_i=\sum\abs{v_{ij}}\epsilon_j$.

\begin{corollary}[Robust verification; \stU]\label{li:cor:robust}
Valid enclosures give $0\le R_K\le R^+$, $\abs{N_i-\widehat S_i}\le E_i+C_iR^+$ and
$N_{\mathrm{miss}}\le R^+/(\tau^K\prod_ir_i^2)$. A sufficient rounding test for a proposed integer
$n_i\ge0$ is $\abs{\widehat S_i-n_i}+E_i+C_iR^+<\frac12$. All operations are rational for rational
data.
\end{corollary}

If $R^+<0$ the enclosures contradict the nonnegative residual and must be rejected.

\begin{theorem}[Eventual certificates; \stU]\label{li:thm:eventual}
Let $\tau>0$ lie outside the spectrum and let $m$ be the number of distinct spectral values above
$\tau$. There is a certificate of the preceding form using $2m+1$ consecutive moments, with
rational centers and radii in arbitrarily narrow disjoint intervals about those values. If $\tau$ is
rational and the moments can be enclosed to any precision, an exhaustive rational search
eventually finds a complete correct certificate above $\tau$.
\end{theorem}

\begin{proof}
Let $m\ge1$ and take the centers at the actual values, with narrow intervals and $b_i>\tau$. The
other occurrences lie in $(0,\sigma]$ for some $\sigma<\tau$; let $S_{\mathrm{tail}}$ be their first
moment. Then $R_K\le S_{\mathrm{tail}}\sigma^{K-1}\prod_ia_i^2$, while $C_i=O((K+1)^2b_i^{-K})$, so
$R_K/(\tau^K\prod r_i^2)\to0$ and $C_iR_K\to0$; the rounding test holds for large $K$, since
$\abs{S_i-N_i}+C_iR_K\le2C_iR_K$. Fix such a $K$ with strict slack. The formulas are continuous in
centers, radii and moments, so nearby rational centers and radii enclosing the same values, and
accurate rational enclosures, keep all strict inequalities. For $m=0$ the test $p_K<\tau^K$
eventually holds since $p_K\le p_1\sigma^{K-1}$. Dovetailing over rational configurations, $K$ and
precisions terminates.
\end{proof}

If the leading values $\lambda_i\le B$ have total multiplicity $M$, $\abs{a_i-\lambda_i}\le\eta$ and the
tail lies below $\sigma<\min a_i$, then
$R_K\le MB^K\eta^2(2B+\eta)^{2m-2}+S_{\mathrm{tail}}\sigma^{K-1}\prod_ia_i^2$: larger $K$ may demand more
accurate centers, and close centers enlarge $D_i^{-2}$.

\begin{corollary}[Complete zeta-zero cluster certificates; \stG]\label{li:cor:zetacluster}
Take $x_n=(\frac14+\gamma_n^2)^{-1}$ over the positive ordinates with multiplicity, with moments from
\cref{li:prop:moments}. If $0<\tau<a_i-r_i<a_i+r_i<4$ and the tests of \cref{li:cor:robust} succeed,
then every positive ordinate $\gamma\le T=\sqrt{\tau^{-1}-\frac14}$ lies in one of the disjoint
intervals
$J_i=\bigl(\sqrt{(a_i+r_i)^{-1}-\frac14},\sqrt{(a_i-r_i)^{-1}-\frac14}\bigr)$, which contains exactly
$n_i$ zeros with multiplicity. If $\tau$ is not a spectral value, arbitrarily narrow such
certificates exist.
\end{corollary}

\begin{proof}
The map $x\mapsto\sqrt{x^{-1}-1/4}$ is decreasing on $(0,4)$; apply
\cref{li:thm:residual,li:thm:eventual}.
\end{proof}

For a primitive nonprincipal $\chi$ the product of the completions of $\chi$ and $\bar\chi$ is even
about $\frac12$; under \cref{src:thm:dgrh} its noncentral zeros lie on the line, and after removing
the central factor the even paired product certifies combined noncentral counts for the pair. For a
real character one may instead use the single completion after removing a central zero. Central
vanishing is not excluded by Dirichlet GRH and must be treated separately.

\begin{example}[An exact certificate with an infinite tail]\label{li:ex:tail}
The multiset $1,\frac12,\frac12,\frac14,\frac1{16},\frac1{32},\frac1{64},\dots$ has moments
$p_j=1+2\cdot2^{-j}+4^{-j}+16^{-j}/(1-2^{-j})$. With centers $(1,\frac12,\frac14)$, radii
$(\frac1{20},\frac1{40},\frac1{80})$ and $\tau=\frac18$, exact rational arithmetic first succeeds at
$K=25$, using $p_{25},\dots,p_{31}$: the missed-count bound is about $0.722<1$ and the counts are
$(1,2,1)$, with error radii about $9.9\cdot10^{-29}$, $1.1\cdot10^{-19}$ and $4.3\cdot10^{-12}$
(\artifact{grh_new_results_2026-11-26/artifacts/certificate_results.json}). A floating-point zeta run
with $K=16$ suggests one zero in $(13.79,14.50)$ and none other below $16.9$; it is not certified.
\end{example}

Recovery of a nonnegative measure on $s$ points from $2s+1$ moments, and the conditioning of Prony
systems, are classical \cite{DeCastroGamboa2012Exact,BatenkovYomdin2013Accuracy}.

Finite moment data certify local statements but cannot force global simplicity, even within real
zero entire functions with the same eventual zero count.

\begin{theorem}[Finite-jet surgery; \stU]\label{li:thm:surgery}
Let $F(z)=F(0)\prod_{n\ge1}(1-a_nz^2)$ with $F(0)\in\R\setminus\{0\}$, $a_n>0$, $\sum a_n<\infty$ and
infinitely many distinct $a_n$, listed with multiplicity. Fix $d,r\ge1$, $d$ distinct anchor nodes
$a_1,\dots,a_d$ and $\epsilon>0$. There is $L>0$ such that for any $r$ other factor occurrences
with positive zero ordinates $\beta_1,\dots,\beta_r\ge L$ and any $T\ge L$ with $F(T)\ne0$, there are
distinct positive $u_i$ with $\abs{u_i-a_i}<\epsilon$ for which
\begin{equation}\label{li:eq:surgery}
 G(z)=F(z)\frac{\prod_{i=1}^d(1-u_iz^2)(1-z^2/T^2)^r}{\prod_{i=1}^d(1-a_iz^2)\prod_{j=1}^r(1-z^2/\beta_j^2)}
\end{equation}
has the following properties.
\begin{enumerate}[label=\textup{(\roman*)}]
\item $G$ is even, entire, real on the real axis and has only real zeros; its zeros at $\pm T$ have
 multiplicity exactly $r$.
\item $G(0)=F(0)$ and $G-F=O(z^{2d+2})$; the first $d$ nonconstant logarithmic Taylor coefficients
 agree.
\item The counts $N_H(Y)$ of positive zeros up to $Y$ satisfy $\abs{N_G(Y)-N_F(Y)}\le d+r$ for all
 $Y\ge0$ and $N_G(Y)=N_F(Y)$ for large $Y$.
\item $\log M_G(R)-\log M_F(R)=O(1)$ for $M_H(R)=\max_{\abs z=R}\abs{H(z)}$; order and type are
 preserved.
\end{enumerate}
\end{theorem}

\begin{proof}
Put $b_j=\beta_j^{-2}$, $c=T^{-2}$ and $\Phi(u)=(\sum_iu_i^k)_{k=1}^d$, with
$\det D\Phi(a)=d!\prod_{i<j}(a_j-a_i)\ne0$. By the inverse function theorem a neighbourhood of
$\Phi(a)$ has preimages in small disjoint positive intervals about the anchors. We solve
$\sum_iu_i^k=\sum_ia_i^k+\sum_jb_j^k-rc^k$, $1\le k\le d$; for $b_j,c\le\eta\le1$ each perturbation is at
most $r\eta$, so small $\eta$ and $L\ge\eta^{-1/2}$ suffice, and $L$ is increased beyond the anchor
ordinates and their perturbations. The removed factors are factors of $F$, so $G$ is entire, and its
node sequence is positive and summable. Since $F(T)\ne0$ the multiplicity at $T$ is $r$. From
$\log(F(z)/F(0))=-\sum_kz^{2k}k^{-1}\sum_na_n^k$ the chosen equations cancel the first $d$ terms of
$\log(G/F)$. Exactly $d+r$ positive zero occurrences were removed and inserted. Finally
$G/F\to c^r\prod u_i/(\prod a_i\prod b_j)>0$ as $\abs z\to\infty$.
\end{proof}

The inverse function theorem may be replaced by a finite certificate: with
$q_k=\sum_ia_i^k+\sum_jb_j^k-rc^k$ and Newton's identities $e_k=\frac1k\sum_{j=1}^k(-1)^{j-1}e_{k-j}q_j$,
opposite endpoint signs of $X^d-e_1X^{d-1}+\cdots+(-1)^de_d$ on disjoint intervals about the anchors
give all its roots \cite{MelanovaSturmfelsWinter2021Recovery}. The anchors may be chosen above any
fixed height, leaving the zeros below it unchanged, and $\sup_{\abs z\le R}\abs{G-F}$ can be made
arbitrarily small.

\begin{example}[A double zero planted in the sine product]\label{li:ex:surgery}
For $F(t)=\sin(\pi t)/(\pi t)$ use the anchors $1,2$, delete the zeros $10,11$ and plant a double
zero at $T=25/2$. With $q_1=189889/151250$ and $e_2=194206753/756250000$, the polynomial
$X^2-q_1X+e_2$ has discriminant $3139703802/5719140625>0$ and roots in $(99/100,101/100)$ and
$(6/25,27/100)$, so
\[
 G(t)=\frac{\sin(\pi t)}{\pi t}\cdot
 \frac{(1-q_1t^2+e_2t^4)(1-4t^2/625)^2}{(1-t^2)(1-t^2/4)(1-t^2/100)(1-t^2/121)}
\]
is entire with only real zeros, the relocated anchors are about $1.0009018469$ and $1.9715556473$,
the pair $\pm25/2$ is exactly double and every other zero is simple. Its derivatives through order
five agree with those of $F$, and $N_G(Y)=\floor Y$ for $Y\ge25/2$. The rational checks are in
\begin{center}
\artifact{grh_new_results_2026-11-26/artifacts/surgery_results.json}.
\end{center}
\end{example}

Under \cref{src:thm:dgrh} the paired product of $\Xi(t)=\xi(\frac12+it)$ satisfies the hypotheses of
\cref{li:thm:surgery}. Thus any finite central jet and any finite window of exact zeros are
compatible, within real-zero entire functions, with an arbitrarily high multiple zero and the same
eventual counting function. No Euler product or arithmetic explicit formula is preserved, so this
does not decide the multiplicities of $\Xi$. Finite-jet nonrigidity by Hermite interpolation was
recorded by Eberwein \cite{Eberwein2026Truncated}.

\subsection{A universal process from the Cayley angles}\label{li:sec:univ}
\sources{rh_li_universality_2026-11-26}

In this subsection $\liZ$ is a locally finite multiset of real numbers with multiplicities, and
$N(T)=\#\{\gamma\in\liZ:\abs\gamma\le T\}$ satisfies, for constants $d>0$ and $B\in\R$,
\begin{equation}\label{li:eq:count}
 N(T)=\frac d\pi T\log T+\frac B\pi T+O(\log(T+2))\qquad(T\to\infty).
\end{equation}
Put $\theta(\gamma)=2\arctan(1/(2\gamma))$ for $\gamma\ne0$, $\theta(0)=\pi$,
$\nu=\frac12\sum_{\gamma\in\liZ}(\delta_{\theta(\gamma)}+\delta_{-\theta(\gamma)})$, so that $m$ occurrences
of $0$ give mass $m/2$ at each of $\pm\pi$, and
\begin{equation}\label{li:eq:symbol}
 \Psi(u)=\int_\R(1-\cos uy)\,\nu(\dd y)=\sum_{\gamma\in\liZ}(1-\cos(u\theta(\gamma))).
\end{equation}
Euler's constant is $\gamma_E$ here.

\begin{lemma}[The process; \stU]\label{li:lem:levy}
$\nu$ is a symmetric L\'evy measure on $[-\pi,\pi]$ with $\nu(\R)=\infty$, $\int\abs y\,\nu(\dd y)=\infty$
and $0<V:=\int y^2\,\nu(\dd y)<\infty$. There is a symmetric L\'evy process $(X_t)_{t\ge0}$ with
$\liEx e^{iuX_t}=e^{-t\Psi(u)}$; it has no Gaussian part, infinitely many jumps, infinite variation,
and $\liEx X_t^2=tV$.
\end{lemma}

\begin{proof}
By \eqref{li:eq:count} and partial summation $\sum_{\abs\gamma>1}\gamma^{-2}<\infty$ and
$\sum_{\abs\gamma>1}\abs\gamma^{-1}=\infty$, and $\abs{\theta(\gamma)}\sim\abs\gamma^{-1}$. L\'evy--Khintchine
gives the process with zero Gaussian coefficient and drift, the variation criterion gives infinite
variation, and two derivatives of the characteristic function at zero give the variance.
\end{proof}

\begin{lemma}[The symbol at large frequency; \stU]\label{li:lem:symbol}
With $A=d/2$ and $\kappa=\gamma_E+B/d$,
\[
 \Psi(u)=A\abs u(\log\abs u+\kappa)+O(\sqrt{\abs u}\log(\abs u+2))\qquad(\abs u\to\infty).
\]
\end{lemma}

\begin{proof}
Take $u\ge2$ and $T=\sqrt u$. Ordinates with $\abs\gamma\le T$ contribute $O(T\log(T+2))$. Beyond,
$\theta(\gamma)=\gamma^{-1}+O(\abs\gamma^{-3})$, and replacing $\theta$ by $1/\gamma$ costs
$O(u\int_T^\infty v^{-3}\,\dd N(v))=O(\log u)$. Let $M(v)=(dv\log v+Bv)/\pi$ and $R=N-M$. For
$g(v)=1-\cos(u/v)$, integration by parts and $\abs{g'(v)}\le u/v^2$ give
$\int_T^\infty g\,\dd R=O(\log(T+2)+u\int_T^\infty\log(v+2)v^{-2}\,\dd v)=O(\sqrt u\log u)$, and extending the
main integral down to zero costs $O(T\log(T+2))$. Hence
\begin{align*}
 \Psi(u)&=\frac1\pi\int_0^\infty(1-\cos(u/v))(d\log v+d+B)\,\dd v+O(\sqrt u\log u)\\
 &=\frac u\pi[(d\log u+d+B)I_0-dI_1]+O(\sqrt u\log u),
\end{align*}
with $I_0=\int_0^\infty(1-\cos y)y^{-2}\dd y=\pi/2$ and
$I_1=\int_0^\infty(1-\cos y)y^{-2}\log y\,\dd y=\frac\pi2(1-\gamma_E)$. Both follow by differentiating
$\int_0^\infty(1-\cos y)y^{-1-\alpha}\dd y=\pi/(2\Gamma(1+\alpha)\sin(\pi\alpha/2))$, $0<\alpha<2$, at
$\alpha=1$; this identity follows from the Laplace representation of $y^{-1-\alpha}$, Tonelli's theorem
and $\int_0^\infty e^{-ty}(1-\cos y)\dd y=1/[t(1+t^2)]$.
\end{proof}

Let $W_0$ be the positive branch of Lambert's function and put
\begin{equation}\label{li:eq:scale}
 w(t)=W_0\Bigl(\frac{e^\kappa}{At}\Bigr),\qquad a(t)=Atw(t),\qquad\text{so}\quad a=e^{\kappa-w},\quad
 \log(1/a)+\kappa=w,\quad t=\frac a{Aw}.
\end{equation}
For $k\ge0$ define, with the principal logarithm,
\begin{equation}\label{li:eq:profiles}
 Q_k(x)=\frac{(-1)^k}{\pi k!}\Re\Bigl.\frac{\partial^k}{\partial z^k}\bigl[\Gamma(z)(1-ix)^{-z}\bigr]\Bigr|_{z=k+1}.
\end{equation}

\begin{theorem}[Universal expansion at every logarithmic order; \stU]\label{li:thm:allorders}
For every $t>0$, $X_t$ has a density $f_t$ that extends to an entire function. For fixed integers
$K,r\ge0$, as $t\downarrow0$,
\[
 \sup_{x\in\R}\Bigl|\frac{\dd^r}{\dd x^r}\Bigl[a(t)f_t(a(t)x)-\sum_{k=0}^K\frac{Q_k(x)}{w(t)^k}\Bigr]\Bigr|
 =O_{K,r}\bigl(w(t)^{-K-1}+\sqrt{a(t)}\bigr)=O_{K,r}(w(t)^{-K-1}).
\]
The constants depend on the multiset and its counting estimate.
\end{theorem}

\begin{proof}
\Cref{li:lem:symbol} makes $\exp(-t\Psi(u)+c\abs u)$ integrable for all $c>0$, so Fourier inversion
gives an entire density. Put $F_w(s)=s+s\log s/w$ for $s>0$; by \eqref{li:eq:scale},
$tA\frac sa(\log\frac sa+\kappa)=F_w(s)$. We first show
\begin{equation}\label{li:eq:fcomp}
 \int_0^\infty s^r\bigl|e^{-t\Psi(s/a)}-e^{-F_w(s)}\bigr|\,\dd s=O_r(\sqrt a).
\end{equation}
On $0<s<\sqrt a$ both exponentials are bounded (the negative part of $F_w$ is at most $e^{-w-1}/w$).
On $s\ge\sqrt a$ both exponents are at least $cs$ for small $t$, since $F_w(s)/s\ge\frac12+\kappa/(2w)$
and the remainder in \cref{li:lem:symbol} is a small fraction of its main term for
$s/a\ge a^{-1/2}$; the exponential mean value bound gives
$\abs{e^{-t\Psi(s/a)}-e^{-F_w(s)}}\le Ct\sqrt{s/a}\log(2+s/a)e^{-cs}$, and $t=a/(Aw)$ with
$\log(2+s/a)\le w+C_\kappa+\log(2+s)$ gives \eqref{li:eq:fcomp}. Next, Taylor's remainder for
$e^{-z}$, $z\ge0$, and $\abs{s\log s}\le1/e$ on $(0,1)$ give
\[
 \int_0^\infty s^re^{-s}\Bigl|e^{-s\log s/w}-\sum_{k=0}^K\frac{(-s\log s)^k}{k!\,w^k}\Bigr|\,\dd s=O_{K,r}(w^{-K-1}).
\]
These are $L^1$ bounds for the Fourier transforms with polynomial weights; Fourier inversion and
differentiation give the uniform claim, and
$\int_0^\infty s^k(\log s)^ke^{-(1-ix)s}\dd s=\partial_z^k[\Gamma(z)(1-ix)^{-z}]_{z=k+1}$ identifies the
profiles. As $a=e^{\kappa-w}$, $\sqrt a$ is smaller than every power of $1/w$.
\end{proof}

The first profiles are
\begin{gather*}
 Q_0(x)=\frac1{\pi(1+x^2)},\qquad Q_1(x)=-\frac1\pi\Re\frac{1-\gamma_E-\log(1-ix)}{(1-ix)^2},\\
 Q_2(x)=\frac1\pi\Re\frac{(\frac32-\gamma_E-\log(1-ix))^2+\frac{\pi^2}6-\frac54}{(1-ix)^3},
\end{gather*}
so the leading shape is the standard Cauchy density.

\begin{corollary}[Universality beyond every logarithmic order; \stU]\label{li:cor:universal}
If two multisets satisfy \eqref{li:eq:count} with the same $d,B$, their scaled densities differ by
$O_r(\sqrt a)$ in $C_b^r(\R)$, and in particular have the same expansion to every order in
$1/\log(1/t)$.
\end{corollary}

\begin{proof}
Both scaled characteristic functions differ from $e^{-F_w(\abs s)}$ by $O_r(\sqrt a)$ in the weighted
$L^1$ norm \eqref{li:eq:fcomp}; apply Fourier inversion.
\end{proof}

For the ordinates of a fixed primitive Dirichlet $L$-function of conductor $q$, counted with
multiplicity and both signs, $N_\chi(T)=\frac T\pi\log\frac{qT}{2\pi e}+O_\chi(\log(T+2))$
\cite{BennettMartinOBryantRechnitzer2021Counting}, and $q=1$ gives $\zeta$. Thus $d=1$,
$B=\log(q/(2\pi))-1$, and
\[
 w_q(t)=W_0\Bigl(\frac{qe^{\gamma_E-1}}{\pi t}\Bigr),\qquad a_q(t)=\frac t2w_q(t)\sim\tfrac12t\log(1/t).
\]
These density statements are unconditional, and all primitive characters of one conductor share the
expansion to every logarithmic order; no uniformity in $q$ is asserted. By \cref{src:thm:dgrh}
every zero is $\rho=\frac12+i\gamma$ with $1-1/\rho=e^{i\theta(\gamma)}$ (with $e^{i\pi}=-1$ at a central zero),
so the Li coefficients satisfy
\begin{equation}\label{li:eq:liident}
 \Re\lambda_\chi(n)=\sum_\rho\Re\bigl[1-(1-1/\rho)^n\bigr]=\Psi_\chi(n),\qquad
 \liEx e^{inX_t}=e^{-t\Re\lambda_\chi(n)}\quad(\text{\stG}),
\end{equation}
and $\liEx e^{inX_t}=e^{-t\lambda_n}$ for $\zeta$. For nonreal characters the real part is essential.

Reduce $X_t$ modulo $2\pi$, with density $p_t(y)=\frac1{2\pi}\sum_{n\in\Z}e^{-t\Psi(n)}e^{iny}$ on
$[-\pi,\pi]$, an entire periodic function. By \eqref{li:eq:liident} the generator $\mathcal A$ of this
circle semigroup has eigenvalues $\Re\lambda_\chi(\abs n)$.

\begin{proposition}[Circle kernel and heat trace; \stU]\label{li:prop:trace}
For fixed $R>0$ and $K,r\ge0$,
$\sup_{\abs x\le R}\bigl|\frac{\dd^r}{\dd x^r}[ap_t(ax)-\sum_{k\le K}Q_k(x)w^{-k}]\bigr|=O_{R,K,r}(w^{-K-1}+\sqrt a)$,
and
\[
 \Tr(e^{-t\mathcal A})=2\pi p_t(0)=\frac2a\Bigl[1-\frac{1-\gamma_E}w+\frac{(\frac32-\gamma_E)^2+\frac{\pi^2}6-\frac54}{w^2}
 +O(w^{-3})\Bigr].
\]
For a primitive Dirichlet character the leading term is $4/(tw_q(t))$.
\end{proposition}

\begin{proof}
The scaled series has frequency mesh $a$. Splitting at $an=\sqrt a$ and the estimates of
\eqref{li:eq:fcomp} bound its difference from the model Riemann sum by $O_r(\sqrt a)$. The model
function $s^re^{-F_w(s)}e^{ixs}$ has uniformly integrable derivative for $\abs x\le R$, so the
Riemann-sum error is $O_{R,r}(a)$; the Taylor step is unchanged, and $x=0$ gives the trace.
\end{proof}

Uniformity over all real $x$ fails on the circle: periodicity produces returning peaks at distance
$2\pi/a$.

\begin{theorem}[Moment transition; \stU]\label{li:thm:moments}
Let $L=\log(1/t)$. As $t\downarrow0$,
\[
 \liEx\abs{X_t}^p\sim\frac{a(t)^p}{\cos(\pi p/2)}\ (0<p<1),\qquad
 \liEx\abs{X_t}\sim\frac d{2\pi}tL^2,\qquad
 \liEx\abs{X_t}^p\sim t\int\abs y^p\,\nu(\dd y)\ (1<p<2),
\]
and $\liEx X_t^2=tV$. Moreover $X_t/a(t)\Rightarrow\mathrm{Cauchy}(0,1)$ as $t\downarrow0$ and
$X_t/\sqrt{tV}\Rightarrow N(0,1)$ as $t\to\infty$.
\end{theorem}

\begin{proof}
For $0<p<2$, $\liEx\abs{X_t}^p=c_p\int_0^\infty(1-e^{-t\Psi(u)})u^{-p-1}\dd u$ with
$c_p=2\Gamma(1+p)\sin(\pi p/2)/\pi$, by Tonelli and the cosine integral of \cref{li:lem:symbol}. For
$0<p<1$ put $u=s/a$: the range $s\le aU_0$ contributes $O(ta^{-p})=o(1)$ to the normalized moment by
$\Psi(u)\le Vu^2/2$; on $aU_0\le s\le1$, $t\Psi(s/a)\le Cs$, and on $s\ge1$, $1-e^{-t\Psi}\le1$;
dominated convergence gives the limit $c_p\int_0^\infty(1-e^{-s})s^{-p-1}\dd s=1/\cos(\pi p/2)$. For
$1<p<2$ divide by $t$ and use $(1-e^{-t\Psi})/t\le\Psi$ with the majorant from $\Psi=O(u^2)$ at $0$ and
\cref{li:lem:symbol} at infinity. For $p=1$ the range $u<1$ is $O(t)$; splitting at $1/t$ with
$1-e^{-z}\le\min(z,1)$ gives the upper bound $\frac A\pi tL^2+O(tL)$, and keeping
$1\le u\le1/(tL^2)$ with $1-e^{-z}\ge z-z^2/2$ gives the lower bound $\frac A\pi t(L-2\log L)^2+O(tL)$.
The weak limits follow from $t\Psi(s/a)\to\abs s$ and from $\Psi(u)=Vu^2/2+o(u^2)$ at $0$.
\end{proof}

The Cauchy limit has no variance because second moments are not uniformly integrable under the
small-time scaling. Nakamura and Suzuki relate RH to infinitely divisible distributions built from
ordinates with inverse-square weights \cite{NakamuraSuzuki2023Infinitely}; here the jumps are the
transformed angles with equal weights. General small-time moment asymptotics for L\'evy processes are
in \cite{FigueroaLopez2008Small}.

\subsection{Open questions}\label{li:sec:open}

\begin{question}[Euler products and weaker recurrence]\label{li:q:euler}
Can a universal Li sampling test requiring less than simultaneous Bohr recurrence be proved for a
class of functions restricted by an Euler product? The perturbations of \cref{li:cor:xi} preserve
the functional equation and the high zeros, whereas an Euler product constrains all Dirichlet
coefficients.
\end{question}

\begin{question}[Frequencies of sign patterns]\label{li:q:signs}
What are the limiting frequencies of the four sign regions of $(d_n,d_{n^2})$ in
\cref{li:cor:signs}, away from their boundaries? A quantitative first-return bound for the certified
patterns of \cref{li:thm:curv} would also give a search length for further examples.
\end{question}

\begin{question}[Effective rigidity]\label{li:q:rigidity}
Give an effective window length in \cref{li:thm:rigidity}: for two number fields with different zeta
functions, bound the length $H$ beyond which every window has mean $\abs{D_n}$ above a given size.
This requires quantitative control of ordinates, of their separation and of the tails as the fields
vary.
\end{question}

\begin{question}[Finer Ritz asymptotics]\label{li:q:ritz}
Does the remainder $O_j(N)$ in \eqref{li:eq:ritzrate}, divided by $N$, have a limit? How does the
estimate change when $j$ grows with $N$? The proof uses the isolation of each of the first $j$ atoms;
uniform control would need quantitative information on their spacings and on the tail.
\end{question}

\begin{question}[Economy of the matched-pair test]\label{li:q:economy}
Minimize the largest Taylor index needed in \cref{li:thm:gap} for a prescribed gap bound, by
optimizing the nonnegative trigonometric kernel in place of $K_N^r$. The simplicity of all zeros of
$\zeta$ and the global disjointness of the ordinates of $\zeta$ and $\beta$ remain open.
\end{question}
